\chapter{Graded first-order rigidity of affine \texorpdfstring{$\mathfrak{sl}_2$}{sl(2)}: Linshaw--Qi and admissible levels}\label{ch:L09}

\providecommand{\lninesl}{\mathfrak{sl}_2}
\providecommand{\lnineg}{\mathfrak{g}}
\providecommand{\lnineV}{\mathcal{V}}
\providecommand{\lnineone}{\mathbf{1}}
\providecommand{\lnineDer}{\operatorname{Der}}
\providecommand{\lninead}{\operatorname{ad}}
\providecommand{\lninehat}{\widehat{\mathfrak{sl}}_2}

This chapter documents the release of the library named in the front matter. The chapter was first written against an earlier commit (the ``documented
commit'' of the historical records below); the mathematics it describes was completed in the commit
immediately before that one (``Prove uniform negative-admissible affine sl2 graded
rigidity''). Quotations from the repository are printed in \texttt{leancode} boxes, which
\texttt{scripts/textbook/sync.py} checks to occur verbatim in the source files of the release
(docstrings and comments are omitted from the quotations, never the code);
tool output is printed in plain boxes.

The question of this chapter is easy to state. Take the simple affine vertex operator algebra
$L_k(\lninesl)$ --- the chiral algebra of the $SU(2)$ Wess--Zumino--Witten model at level $k$ --- and
ask whether its operator product expansion can be deformed. Not deformed by a change of
coordinates (renaming the fields, which changes nothing), but genuinely: is there a one-parameter
family of vertex algebras $V_t$, with $V_0 = L_k(\lninesl)$, that is not just $L_k(\lninesl)$ in
disguise? A VOA with no such infinitesimal family is called \emph{rigid}. The library works to first
order and with gradings: its notion is \emph{graded first-order rigidity} (deformations over
$\C[t]/(t^2)$, the vertex operation deformed, the vacuum fixed, the weight grading respected; a
deformation is trivial when it is the coboundary of a grade-preserving, vacuum-fixing $\varphi$). Every
rigidity statement in this chapter is meant in that sense, and nothing is claimed about formal,
ungraded or grading-moving deformations. The library proves:

\begin{theorem}[informal; precise form in Theorem~\ref{l09:thm:headline}]
Let $p,q$ be coprime positive integers with $2 \le p < 2q$, and put $k = p/q - 2$. Then every
graded first-order deformation of the simple affine VOA $L_k(\lninesl)$ is trivial.
\end{theorem}
\begin{physics}
The $SU(2)$ WZW model at level $k$ is the CFT of a string moving on the group manifold $SU(2)\cong S^3$
with $k$ units of $H$-flux; its holomorphic currents $J^a(z)$ generate $L_k(\lninesl)$. At negative fractional
$k$ the corresponding CFTs are non-unitary (negative central charge $c = 3k/(k+2)$, e.g.\ $c=-6$ at $k=-4/3$).
The theorem says such a chiral algebra admits no infinitesimal change of its OPEs other than field
redefinitions: there is no ``chiral conformal manifold'' through it. A physicist expects this because the
only continuous parameter of the current OPE \eqref{l09:eq:currentOPE} is $k$, and at admissible $k$ the
null-vector relations exist only at that value. This is a heuristic; the proof is below.
\end{physics}


The levels covered are exactly the \emph{negative admissible} levels of Kac and Wakimoto:
$k = -4/3$, $-1/2$, $-8/5$, $-2/3$, $-12/7$, $-5/4$, \ldots{} (every rational number $k\in(-2,0)$
whose shifted level $k+2$ has numerator at least $2$ in lowest terms). The case $k=-4/3$ is
Theorem~5.1 of Linshaw and Qi (arXiv:2601.12017); the uniform statement is new in this library.

\begin{status}
The headline is \lean{VOA.AffineSl2GeneralRigidity.negative_admissible_graded_rigidity} in
\leanpath{VOA/Examples/AffineSl2AdmissibleRigidity.lean}. Its only hypotheses are the four numerical
conditions $2\le p$, $0<q$, $\gcd(p,q)=1$, $p<2q$; no singular-vector, determinant, character,
representation-theoretic or reconstruction premise remains. The Linshaw--Qi level is separately
\lean{VOA.LQ43.L43_graded_rigidity} (no hypotheses at all). While writing this chapter,
\texttt{\#print axioms} was re-run on both (with \texttt{lake env lean -DmaxSynthPendingDepth=3}
against the built tree), and on every other headline declaration of the chapter; the answer was in
every case \texttt{[propext, Classical.choice, Quot.sound]}. The PBW theorem \lean{VOA.pbw} (Chapter~\ref{ch:L02}; formerly a project axiom, now a theorem) is
\emph{not} used anywhere in this chapter's results, even though they concern an affine
Lie-algebra construction: the affine vacuum module is shown nonzero through a PBW theorem
\emph{proved} by Wakimoto triangularity, not through \lean{VOA.pbw}. Both theorems are therefore
\emph{unconditional} statements of graded first-order rigidity. Two caveats from the release audit
(\texttt{Notes/Textbook/referee/RELEASE\_AUDIT\_RIGIDITY.md}) accompany them: only graded first-order
deformations are covered, and the non-vacuity of the rigidity notion (that some VOA fails it) was argued on
paper only. The second caveat has since been discharged: the library proves that the universal Virasoro VOA
$V(c,0)$ is not graded rigid, for every $c$ (\lean{VOA.}\allowbreak\lean{VirPoly.}\allowbreak\lean{not_}\allowbreak\lean{isRigidGraded_}\allowbreak\lean{virasoroVOA}; the witness $\partial_cY$ moves the
central charge, which \lean{FirstOrderDeformation} permits; Section~\ref{l09:sec:deform}, second warning). The full record is in
Section~\ref{l09:sec:verification}.
\end{status}

\paragraph{How to read this chapter.}
Every definition and theorem appears as a mathematical statement, then its Lean declaration
copied verbatim, then a fidelity-check box (titled \emph{From the mathematics to the Lean
statement}) that decodes the Lean, lists every hypothesis, traces every defined term to its
definition, compares with the literature, and says how the statement could have been wrong or
vacuous. Proofs are explained on paper. Sections~\ref{l09:sec:affine}--\ref{l09:sec:deform} set up
the objects (affine $\lninesl$, admissible levels, deformations); Section~\ref{l09:sec:whitehead} is
the Lie-algebra cohomology needed to normalise a deformation;
Section~\ref{l09:sec:lq} is the story of Linshaw--Qi's Theorem~5.1, including the
\emph{tautological bundled hypothesis} that an earlier revision of the library carried and how it
was removed; Sections~\ref{l09:sec:uniform}--\ref{l09:sec:headline} give the uniform proof;
Section~\ref{l09:sec:boundary} describes the (conditional) boundary-level $C_2$ programme that lives
in the same files. The chapter ends with the verification record, the referee's checklist, and
notes and references.

\paragraph{A physicist's preview.}
In a two-dimensional conformal field theory, the chiral algebra is the algebra of holomorphic
fields and their operator product expansions (OPEs). For the $SU(2)$ WZW model the chiral algebra
is generated by three currents $J^e(z), J^h(z), J^f(z)$ with OPE
\begin{equation}\label{l09:eq:currentOPE}
  J^a(z)\,J^b(w) \sim \frac{k\,\kappa(a,b)}{(z-w)^2} + \frac{J^{[a,b]}(w)}{z-w},
\end{equation}
where $\kappa$ is the normalised Killing form. The OPE has exactly two pieces of data: the
structure constants of $\lninesl$ (which cannot change continuously --- $\lninesl$ is rigid as a Lie
algebra) and the number $k$ in the double pole. A first-order deformation of the OPE therefore
has, morally, one candidate direction: \emph{change the level}, $k \mapsto k + \delta t$. For the
\emph{universal} affine algebra $V^k(\lninesl)$ this is a genuine deformation --- the algebras
$V^{k+t}(\lninesl)$ form an honest family. For the \emph{simple} quotient $L_k(\lninesl)$ at an
admissible level something new happens: $L_k$ is defined by imposing a null-vector relation
$s = 0$, and the null vector exists only at that particular $k$. Changing the level breaks the
null relation, so the deformation is inconsistent unless $\delta = 0$. The whole chapter is the
rigorous version of this sentence: the pole parameter $\delta$ is shown to be well defined (after a
gauge fixing that uses Whitehead's lemma), the null relations are shown to force
$\delta \cdot P'(k) = 0$ for an explicit polynomial $P$, and the Kac determinant is used to show
$P'(k)\neq 0$.

\section{Affine \texorpdfstring{$\mathfrak{sl}_2$}{sl(2)} and its simple quotients}\label{l09:sec:affine}

This section recalls, with the Lean source, the objects in the headline theorem. The detailed
construction of the affine VOA (reconstruction, Sugawara vector, Wakimoto realisation, PBW
theorem) is Chapter~\ref{ch:L08}; the general notions of vertex operator algebra, ideals and
simplicity are in Chapters~\ref{ch:L03} and~\ref{ch:L04}. Here we give what a reader needs to parse every statement.

\subsection{The Lie algebra \texorpdfstring{$\mathfrak{sl}_2$}{sl(2)} and its normalised Killing form}

$\lninesl$ is the three-dimensional complex Lie algebra with basis $e,f,h$ and brackets
\begin{equation*}
  [h,e] = 2e,\qquad [h,f] = -2f,\qquad [e,f] = h .
\end{equation*}
Its invariant symmetric bilinear forms form a one-dimensional space; the \emph{normalised Killing
form} $\kappa$ is the one with $\kappa(e,f) = \kappa(f,e) = 1$, $\kappa(h,h) = 2$ and all other basis
pairings zero (it is $\tfrac14$ of the Killing form $\operatorname{tr}(\lninead x\,\lninead y)$, and
equals the trace form $\operatorname{tr}(xy)$ in the defining $2\times 2$ representation). The
library realises $\lninesl$ as coordinate vectors, with the colour convention $e=0$, $f=1$, $h=2$:

\begin{leancode}
abbrev Sl2 : Type := Fin 3 → ℂ
\end{leancode}
\begin{leancode}
def kappa : Fin 3 → Fin 3 → ℂ
  | 0, 1 => 1
  | 1, 0 => 1
  | 2, 2 => 2
  | _, _ => 0
\end{leancode}
(\leanpath{VOA/Examples/AffineSl2.lean}; the bracket is defined there by the structure constants
$f^{ab}{}_c$ above, and \lean{kappaForm} is the bilinear form with matrix \lean{kappa}, proved
invariant as \lean{kappaForm_isInvariant}.)

\subsection{The affine Lie algebra and the vacuum module}

The affine Lie algebra $\lninehat = \lninesl\otimes\C[t,t^{-1}] \oplus \C C \oplus \C\ell_0$ has
brackets
\begin{equation}\label{l09:eq:affinebracket}
  [x\otimes t^m, y\otimes t^n] = [x,y]\otimes t^{m+n} + m\,\kappa(x,y)\,\delta_{m+n,0}\,C,
  \qquad C \text{ central},
\end{equation}
and $\ell_0$ the degree derivation. We write $J^a_m$ for the action of $a\otimes t^m$. Fix a complex
number $k$, the \emph{level}. The \emph{vacuum module} (also called the generalised Verma module
induced from the trivial $\lninesl$-module) is
\begin{equation*}
  V^k(\lninesl) = U(\lninehat)\big/ J(k),\qquad
  J(k) = \text{left ideal generated by } x\otimes t^n\ (n\ge 0),\ C - k,\ \ell_0 .
\end{equation*}
Its vacuum $\lnineone$ is the class of $1$; it is annihilated by all non-negative modes, and $C$
acts as $k$. By the PBW theorem, $V^k(\lninesl)$ has basis the ordered monomials
$J^{a_1}_{-n_1}\cdots J^{a_r}_{-n_r}\lnineone$ with all $n_i\ge 1$; the grading by $L_0$ (energy) is
$\sum n_i$. The library builds exactly this quotient (\leanpath{VOA/Foundations/AffineLie.lean}):

\begin{leancode}
def vacuumGens (k : ℂ) : Set (UEnv 𝔤 κ hκ) :=
  {a | (∃ (x : 𝔤) (n : ℤ), 0 ≤ n ∧ a = ι ℂ (ofLoop κ hκ x n))
        ∨ a = ι ℂ (C κ hκ) - algebraMap ℂ (UEnv 𝔤 κ hκ) k
        ∨ a = ι ℂ (ell0 κ hκ)}
\end{leancode}
\begin{leancode}
noncomputable def vacuumIdeal (k : ℂ) : Submodule (UEnv 𝔤 κ hκ) (UEnv 𝔤 κ hκ) :=
  Submodule.span _ (vacuumGens κ hκ k)
\end{leancode}
\begin{leancode}
abbrev vacuumSpace (k : ℂ) : Type _ := UEnv 𝔤 κ hκ ⧸ vacuumIdeal κ hκ k
\end{leancode}
and, for $\lninesl$, gives the carrier a (historically chosen, slightly misleading) name:
\begin{leancode}
noncomputable abbrev WakModule (k : ℂ) : Type :=
  vacuumSpace kappaForm kappaForm_isInvariant k
\end{leancode}

\begin{dictionary}
\lean{UEnv 𝔤 κ hκ} is Mathlib's universal enveloping algebra of the affine Lie algebra
\lean{AffineLie 𝔤 κ hκ}; \lean{ι ℂ} is the canonical map $\lnineg\to U(\lnineg)$;
\lean{ofLoop κ hκ x n} is $x\otimes t^n$; \lean{C κ hκ} is the central element and \lean{ell0} the
degree derivation. \lean{Submodule.span _ S} with the ring $U$ acting on itself on the left is the
\emph{left} ideal generated by $S$, and \lean{⧸} is the quotient module. So \lean{vacuumSpace k} is
literally $U(\lninehat)/J(k)$ as above. The name \lean{WakModule} is a historical artefact: the
type is the universal vacuum module $V^k(\lninesl)$, \emph{not} a Wakimoto (free-field) module. The
Wakimoto realisation enters only as a tool for proving that this module is nonzero and has a PBW
basis.
\end{dictionary}

\begin{warning}
\textbf{Where \lean{VOA.pbw} could have entered, and why it does not.} A quotient $U/J$ by a left
ideal can be zero (if $1\in J$); a VOA structure on the zero module satisfies every axiom, so
every theorem about it would be vacuously true. The library's PBW theorem
\lean{VOA.pbw} (Chapter~\ref{ch:L02}; an axiom in earlier versions of the library) asserts precisely that affine vacuum vectors are nonzero, and
one lemma in \leanpath{VOA/Examples/AffineSl2VOAInstance.lean} (\lean{closed_vacuum_ne_zero}) uses it.
The results of this chapter avoid that lemma: they take nonvanishing from
\lean{AffineSl2BoundaryMaximalIdeal.affineSl2_vacuum_ne_zero}, which is proved from the affine PBW
basis \lean{AffineSl2WakimotoLiLeading.affineSl2PBWBasis}, itself proved (Chapter~\ref{ch:L08}) by
leading-term triangularity in the Wakimoto realisation. The \texttt{\#print axioms} record of
Section~\ref{l09:sec:verification} confirms that no result here reaches \lean{VOA.pbw}.
\end{warning}

\subsection{The universal affine VOA and its simple quotient}

For $k\neq -2$ (the \emph{critical} level is $k = -h^\vee = -2$) the Sugawara vector
\begin{equation*}
  \omega = \frac{1}{2(k+2)}\Bigl(\tfrac12 h_{-1}h_{-1} + e_{-1}f_{-1} + f_{-1}e_{-1}\Bigr)\lnineone
\end{equation*}
makes $V^k(\lninesl)$ a vertex operator algebra of central charge $c = 3k/(k+2)$, with state--field
map determined by $Y(J^a_{-1}\lnineone, z) = \sum_n J^a_n z^{-n-1}$. The library's closed instance is
\lean{AffineVOA.affineSl2VOA_closed k hk} (\leanpath{VOA/Examples/AffineSl2Virasoro.lean}, built in
Chapter~\ref{ch:L08}); this chapter abbreviates it:

\begin{leancode}
noncomputable abbrev universal := affineSl2VOA_closed k hk
\end{leancode}
\begin{leancode}
noncomputable abbrev simple := QuotVOA.maxQuotientVOA (universal k hk)
\end{leancode}
(\leanpath{VOA/Examples/AffineSl2GeneralRigidityBasis.lean}, under \lean{variable (k : ℂ) (hk : k + 2 ≠ 0)}.)
Here the simple quotient is formed by the general construction of Chapter~\ref{ch:L04}:

\begin{leancode}
def maxGradedIdeal (W : VertexOperatorAlgebra 𝒱) : Submodule ℂ 𝒱 :=
  sSup (ProperGradedIdeals W)
\end{leancode}
\begin{leancode}
def maxQuotientVOA (W : VertexOperatorAlgebra 𝒱) :
    VertexOperatorAlgebra (𝒱 ⧸ maxGradedIdeal W) :=
  quotientVOA W (maxGradedIdeal W) (maxGradedIdeal_isIdeal W) (maxGradedIdeal_isGradedIdeal W)
\end{leancode}
(both in \leanpath{VOA/Core/QuotientVOA.lean}), with ideals and simplicity defined by
(\lean{IsIdeal}, \lean{IsSimple} in \leanpath{VOA/Core/BilinearForm.lean}; \lean{IsGradedIdeal} in
\leanpath{VOA/Core/QuotientVOA.lean}):

\begin{leancode}
def IsIdeal (W : VertexOperatorAlgebra 𝒱) (I : Submodule ℂ 𝒱) : Prop :=
  ∀ (u v : 𝒱) (n : ℤ), v ∈ I → W.smul u n v ∈ I
\end{leancode}
\begin{leancode}
def IsGradedIdeal (W : VertexOperatorAlgebra 𝒱) (I : Submodule ℂ 𝒱) : Prop :=
  I = ⨆ n : ℤ, I ⊓ W.grading n
\end{leancode}
\begin{leancode}
def IsSimple (W : VertexOperatorAlgebra 𝒱) : Prop :=
  (⊤ : Submodule ℂ 𝒱) ≠ ⊥ ∧ ∀ I : Submodule ℂ 𝒱, W.IsIdeal I → I = ⊥ ∨ I = ⊤
\end{leancode}

\begin{definition}\label{l09:def:Lk}
For $k\neq-2$, $L_k(\lninesl) := V^k(\lninesl)/\mathcal{I}_k$, where $\mathcal{I}_k$ is the sum of all
proper graded ideals of $V^k(\lninesl)$ (the maximal proper graded ideal). In Lean,
$L_k(\lninesl)$ is \lean{AffineSl2GeneralRigidity.simple k hk}.
\end{definition}
\begin{physics}
$V^k(\lninesl)$ is the chiral algebra generated freely by the currents subject only to the OPE
\eqref{l09:eq:currentOPE}; $L_k(\lninesl)$ is obtained by setting all null states (states orthogonal to
everything under the BPZ/Shapovalov pairing) to zero, which is what a physical CFT with a nondegenerate
two-point function does. The Sugawara vector $\omega$ is the stress tensor $T(z) = \frac{1}{2(k+2)}(J J)(z)$
built from the currents; $V_1$ is the space of conserved currents, here the $SU(2)$ symmetry currents.
\end{physics}


\begin{dictionary}
\lean{W.smul u n v} is the $n$-th product $u_{(n)}v$, i.e.\ the coefficient of $z^{-n-1}$ in
$Y(u,z)v$ (Chapter~\ref{ch:L03}). An \lean{IsIdeal} is a subspace closed under left multiplication by
all modes of all states (for a VOA, skew-symmetry makes left ideals two-sided). \lean{IsGradedIdeal I}
says $I$ is the sum of its homogeneous pieces. \lean{ProperGradedIdeals W} (in
\leanpath{VOA/Core/QuotientVOA.lean}) is the set of graded ideals not containing the vacuum; its
supremum \lean{sSup} is again a proper graded ideal (a graded ideal is proper iff its weight-zero
part, which is inside $\C\lnineone$, is zero, and that property is preserved under sums). So
\lean{maxGradedIdeal} is the maximal proper graded ideal and \lean{maxQuotientVOA} the quotient VOA,
as in the textbook. \lean{IsSimple} quantifies over \emph{all} ideals, graded or not, and requires
the algebra to be nonzero. The library proves, for every non-critical $k$:
\end{dictionary}

\begin{leancode}
theorem simple_vacuum_ne_zero : (simple k hk).vacuum ≠ 0 :=
  SimpleQuot.maxQuotientVOA_vacuum_ne_zero (universal k hk)
    (AffineSl2BoundaryMaximalIdeal.affineSl2_vacuum_ne_zero k hk)
\end{leancode}
\begin{leancode}
theorem simple_isSimple : (simple k hk).IsSimple :=
  QuotVOA.isSimple_maxQuotientVOA' _
    (AffineSl2BoundaryMaximalIdeal.affineSl2_vacuum_ne_zero k hk)
\end{leancode}

\begin{remark}
The library does not need, and does not prove, a presentation of $L_k(\lninesl)$ by generators
and relations (for example, that $\mathcal{I}_k$ is generated by a single singular vector). The
quotient is defined intrinsically, which is the safest possible choice for a rigidity theorem: the
object is the standard $L_k(\lninesl)$ by definition, not ``$V^k$ modulo whatever relation we wrote
down''. The cost is that facts about $\mathcal{I}_k$ have to be extracted from its definition; the
main tool for that is Li's theorem that $\mathcal{I}_k$ equals the radical of the invariant bilinear
form (Chapter~\ref{ch:L04}), used repeatedly below.
\end{remark}

\subsection{The weight-one space: the currents survive}

The weight-one space of $L_k(\lninesl)$ is spanned by the images of the three currents
$J^a_{-1}\lnineone$. For $k \neq 0$ they remain linearly independent: if
$\sum_a c_a J^a_{-1}\lnineone \in \mathcal{I}_k$, apply the positive modes $J^b_1$, which map into
weight zero; since $J^b_1 J^a_{-1}\lnineone = k\,\kappa(b,a)\lnineone$ and $\lnineone\notin\mathcal{I}_k$,
nondegeneracy of $\kappa$ and $k\neq0$ force all $c_a=0$. (At $k=0$ the opposite happens: the
currents are null, $L_0(\lninesl)=\C$, and $V_1 = 0$.) The library records the result as an
$\lninesl$-triple basis of the weight-one Lie algebra:

\noindent File \leanpath{VOA/Examples/AffineSl2GeneralRigidityBasis.lean}.
\begin{leancode}
theorem simple_sl2BasisOnWeightOne (hk0 : k ≠ 0) : Sl2BasisOnWeightOne (simple k hk) :=
  ⟨(simple_weightOneBasis k hk hk0), (simple_weightOneTriple k hk hk0)⟩
\end{leancode}
where (\leanpath{VOA/Examples/AffineSl2Rigidity43.lean})
\begin{leancode}
def Sl2BasisOnWeightOne {𝒱 : Type*} [AddCommGroup 𝒱] [Module ℂ 𝒱]
    (W : VertexOperatorAlgebra 𝒱) : Prop :=
  ∃ b : Module.Basis (Fin 3) ℂ ↥(W.grading 1), IsSl2Triple (b 1) (b 0) (b 2)
\end{leancode}

\begin{dictionary}
\lean{W.grading 1} is the weight-one subspace $V_1$ (Chapter~\ref{ch:L03}); \lean{↥} turns the
submodule into a type. \lean{Module.Basis (Fin 3) ℂ ↥(W.grading 1)} is a basis of $V_1$ indexed by
$\{0,1,2\}$, and Mathlib's \lean{IsSl2Triple h e f} says $[e,f]=h$, $[h,e]=2e$, $[h,f]=-2f$ and
$h\ne0$, for the Lie bracket on $V_1$ given by $[a,b] = a_{(0)}b$ (Chapter~\ref{ch:L04}). Note the
argument order: \lean{IsSl2Triple (b 1) (b 0) (b 2)} means $b_0 = e$, $b_1 = h$, $b_2 = f$. So
\lean{Sl2BasisOnWeightOne W} is the statement ``$V_1\cong\lninesl$ as Lie algebras''. The theorem
\lean{simple_sl2BasisOnWeightOne} needs $k\ne0$, and this is the only reason the headline theorem
excludes $k=0$.
\end{dictionary}

\section{Admissible levels}\label{l09:sec:admissible}

\subsection{From integrable to admissible}

At a non-negative integer level $k$, $L_k(\lninesl)$ is the chiral algebra of the unitary $SU(2)_k$
WZW model: it is rational and $C_2$-cofinite, it has exactly $k+1$ irreducible modules (spins
$0,\tfrac12,\ldots,\tfrac k2$), and its maximal ideal is generated by the single singular vector
$(e_{-1})^{k+1}\lnineone$. Kac and Wakimoto (1988) discovered that a much larger set of levels
shares the most striking feature of the integrable case --- a finite set of highest-weight modules
whose characters span a representation of $SL_2(\Z)$ --- although the resulting theories are not
unitary.

\begin{definition}[Kac--Wakimoto]\label{l09:def:admissible}
A level $k\in\C$ is \emph{admissible} for $\lninesl$ if
\begin{equation*}
  k + 2 = \frac{p}{q},\qquad p,q\in\Z_{>0},\qquad \gcd(p,q) = 1,\qquad p\ge 2 .
\end{equation*}
It is \emph{negative admissible} if moreover $k<0$, i.e.\ $p<2q$; it is \emph{boundary admissible} if
$p = 2$ (the smallest allowed numerator, which forces $q$ odd).
\end{definition}
\begin{physics}
Admissible levels are the fractional levels at which the $\widehat{\mathfrak{sl}}_2$ characters still
form a finite-dimensional representation of the modular group $SL_2(\Z)$ (Kac--Wakimoto), so they define
non-unitary CFTs with a modular-invariant set of characters: fractional-level $SL(2)$ WZW models, appearing in
Polyakov's 2d gravity, in the Drinfeld--Sokolov/coset description of the $(p,q)$ Virasoro minimal models,
and as chiral algebras of 4d $\mathcal{N}=2$ Argyres--Douglas theories. Negative admissible means negative
level, where the theory is necessarily non-unitary.
\end{physics}


The shift by $2$ is the dual Coxeter number $h^\vee$ of $\lninesl$; for a general simple
$\lnineg$ the condition reads $k+h^\vee = p/q$ with $p\ge h^\vee$ (when $q$ is prime to the lacing
number). The integrable levels are the case $q = 1$, $p = k+2\ge2$. The bound $p\ge2$ is not
decoration: $k = -3/2$ ($k+2 = 1/2$) is not admissible.

\begin{leancode}
def AdmissibleLevel (k : ℂ) : Prop :=
  ∃ p q : ℕ, 2 ≤ p ∧ 1 ≤ q ∧ Nat.Coprime p q ∧ k + 2 = (p : ℂ) / (q : ℂ)
\end{leancode}
(\leanpath{VOA/Examples/AdmissibleLevel.lean}, namespace \lean{VOA.Admissible}.)

\begin{dictionary}
\lean{∃ p q : ℕ} quantifies over natural numbers; \lean{Nat.Coprime p q} is $\gcd(p,q)=1$; the casts
\lean{(p : ℂ)} embed $\N\subset\C$, and \lean{/} is division in $\C$. This is exactly
Definition~\ref{l09:def:admissible} with $q\ge1$ (the inclusive convention that makes integrable
levels admissible; some papers impose $q\ge2$ to talk only about fractional levels, and the file
docstring records that a future fractional-only notion must say \lean{1 < q} explicitly). The
witnesses are unique (\lean{AdmissibleLevel.witness_unique}), so ``the'' $p$ and $q$ of an
admissible level are well defined. \emph{Note for the referee:} the headline theorem does
\emph{not} use this predicate; it takes $p$ and $q$ as explicit natural numbers with the same
conditions plus $p<2q$ (Section~\ref{l09:sec:headline}). The predicate is used for the examples
and for the arithmetic lemmas below.
\end{dictionary}

The library records the elementary consequences:

\begin{leancode}
theorem AdmissibleLevel.add_two_ne_zero {k : ℂ} (h : AdmissibleLevel k) : k + 2 ≠ 0 := by
\end{leancode}
\begin{leancode}
theorem admissibleLevel_boundary (m : ℕ) :
    AdmissibleLevel (-2 + 2 / ((2 * m + 1 : ℕ) : ℂ)) := by
\end{leancode}
\begin{leancode}
theorem centralChargeA_of_add_two_eq {k p q : ℂ} (hp : p ≠ 0) (hq : q ≠ 0)
    (hk : k + 2 = p / q) : centralChargeA k = 3 - 6 * q / p := by
\end{leancode}
and the named examples \lean{admissibleLevel_neg_four_thirds} ($(p,q) = (2,3)$),
\lean{admissibleLevel_neg_half} ($(3,2)$), \lean{admissibleLevel_neg_eight_fifths} ($(2,5)$),
\lean{admissibleLevel_natCast} ($k = n$, $(p,q) = (n+2,1)$). All are pure arithmetic, proved by
\lean{norm_num}/\lean{omega}.

\begin{example}\label{l09:ex:levels}
The first negative admissible levels, ordered by denominator, with central charge
$c = 3k/(k+2) = 3 - 6q/p$:
\begin{center}
\begin{tabular}{cccccc}
\toprule
$q$ & $p$ & $k = p/q-2$ & $c$ & boundary? & weight of first singular vector $q(p-1)$\\
\midrule
$2$ & $3$ & $-1/2$ & $-1$ & no & $4$\\
$3$ & $2$ & $-4/3$ & $-6$ & yes & $3$\\
$3$ & $4$ & $-2/3$ & $-3/2$ & no & $9$\\
$3$ & $5$ & $-1/3$ & $-3/5$ & no & $12$\\
$4$ & $3$ & $-5/4$ & $-5$ & no & $8$\\
$4$ & $5$ & $-3/4$ & $-9/5$ & no & $16$\\
$4$ & $7$ & $-1/4$ & $-3/7$ & no & $24$\\
$5$ & $2$ & $-8/5$ & $-12$ & yes & $5$\\
\bottomrule
\end{tabular}
\end{center}
(For $q = 1$ the condition $2\le p<2$ is empty: integrable levels are never negative.) At
$k=-1/2$ the algebra $L_{-1/2}(\lninesl)$ is the $\Z_2$-orbifold of the $\beta\gamma$ system; at
$k=-4/3$ its representation theory was constructed by Adamovi\'c (2005), and in
the physics literature it appears as the chiral algebra of the $(A_1,D_3)$ Argyres--Douglas theory;
more generally the boundary levels $k = -4n/(2n+1)$ are associated with the $(A_1,D_{2n+1})$
theories (Creutzig; Beem--Rastelli and collaborators). These identifications are background, not
used below.
\end{example}

\subsection{The vacuum singular vector at an admissible level}

Why do admissible levels make $V^k(\lninesl)$ non-simple? The Kac--Kazhdan criterion says that a
Verma module of highest weight $\hat\lambda$ over $\lninehat$ has a singular vector of weight
$\hat\lambda - n\beta$ whenever $\beta$ is a positive real root and
\begin{equation}\label{l09:eq:KK}
  2(\hat\lambda+\hat\rho,\beta) = n(\beta,\beta),\qquad n\in\Z_{>0} .
\end{equation}
For the vacuum, $\hat\lambda = k\Lambda_0$, and with the standard normalisation
$(\alpha,\alpha) = 2$, $(\delta,\Lambda_0) = 1$, $\hat\rho = 2\Lambda_0 + \alpha/2$, the positive real roots
are $m\delta + \alpha$ ($m\ge0$) and $m\delta - \alpha$ ($m\ge1$), with
\begin{equation*}
  (\hat\lambda+\hat\rho, m\delta+\alpha) = m(k+2) + 1,\qquad
  (\hat\lambda+\hat\rho, m\delta-\alpha) = m(k+2) - 1 .
\end{equation*}
The root $\alpha$ itself ($m=0$) gives $n = 1$ for every $k$: this is the singular vector
$f_0\lnineone$, which is already zero in $V^k$ (it is what distinguishes the vacuum module from the
Verma module). At $k+2 = p/q$, the condition for $\beta = m\delta - \alpha$ is $mp/q - 1 = n$, i.e.\
$mp = q(n+1)$; by coprimality $m = \ell q$ and $n = \ell p - 1$, and the smallest solution is
\begin{equation}\label{l09:eq:firstsing}
  \beta = q\delta - \alpha,\quad n = p-1:\qquad
  \text{a singular vector of energy } q(p-1) \text{ and } h_0\text{-weight } 2(p-1).
\end{equation}
For $\beta = m\delta+\alpha$ the condition $mp + q = nq$ forces $m = \ell q$, $n = \ell p + 1$ and energy
strictly above $q(p-1)$. At $k=-4/3$ this is the weight-three singular vector of $h_0$-weight $2$
found by Adamovi\'c (Section~\ref{l09:sec:lq}); for integrable $k$ ($q=1$) it is $(e_{-1})^{k+1}\lnineone$.
Kac--Wakimoto and Malikov--Feigin--Fuchs give explicit formulas; Adamovi\'c--Milas proved that this
vector generates the maximal ideal. None of these literature facts is \emph{assumed} by the
headline theorem: the uniform proof (Section~\ref{l09:sec:determinant}) works with a determinant
instead of an explicit singular vector. The arithmetic half of \eqref{l09:eq:firstsing} is proved
as stated (\leanpath{VOA/Examples/AffineSl2AdmissibleRootArithmetic.lean}):

\begin{leancode}
theorem first_minus_resonance_unique {p q m r : ℕ} (hp : 2 ≤ p) (hq : 0 < q)
    (hm : 0 < m) (hcop : p.Coprime q) (he : m*p = q*(r+1))
    (hbound : m*r ≤ q*(p-1)) : m=q ∧ r=p-1 := by
\end{leancode}
\begin{leancode}
theorem plus_resonance_energy_gt {p q m r : ℕ} (hp : 2 ≤ p) (hq : 0 < q)
    (hm : 0 < m) (hcop : p.Coprime q) (he : m*p+q=q*r) :
    q*(p-1) < m*r := by
\end{leancode}

\begin{dictionary}
In \lean{first_minus_resonance_unique}, $m$ is the $\delta$-coefficient of the root $m\delta-\alpha$,
$r$ is the multiplicity $n$ of \eqref{l09:eq:KK} (so \lean{he} is $mp = q(n+1)$), and \lean{m*r} is
the energy $m\cdot n$ of the corresponding singular weight; the conclusion is that the only
solution with energy at most $q(p-1)$ is $m = q$, $n = p-1$. The natural-number subtraction
\lean{p-1} is honest because $p\ge2$. \lean{plus_resonance_energy_gt} says the other root family
only resonates above that energy. Both are statements about integers only; the file docstring
stresses that ``this proves arithmetic only, not an affine determinant formula or existence of a
singular state''.
\end{dictionary}

\subsection{Physics: WZW models at fractional level}

The admissible $L_k(\lninesl)$ are the chiral algebras of $SL(2)$ WZW models at fractional level,
which appeared in Polyakov's light-cone two-dimensional gravity ($SL(2,\mathbb{R})$ current algebra
at level $k$ with $k+2$ rational), in the coset and Drinfeld--Sokolov constructions of Virasoro
minimal models (the quantum Hamiltonian reduction of $L_k(\lninesl)$ at $k+2 = p/q$ is the $(p,q)$
minimal model, with $c_{\mathrm{Vir}} = 1 - 6(p-q)^2/(pq)$; the library has the arithmetic as
\lean{VOA.DSsl2.dsCentralCharge_eq_minimalModel}), and more recently as chiral algebras of
four-dimensional $\mathcal{N}=2$ superconformal theories. They are not rational in the usual sense
(they have infinitely many irreducible weak modules), not $C_2$-cofinite, but \emph{quasi-lisse}:
the associated variety is the nilpotent cone of $\lninesl$ (Arakawa; a literature fact, not formalised in the
library: associated varieties and quasi-lisseness are not defined there, see Chapter~\ref{ch:L16}). Rigidity
is then a non-trivial statement precisely because none of the usual finiteness properties is
available: Linshaw and Qi's abstract emphasises that their $-4/3$ result shows ``neither
$C_2$-cofiniteness nor rationality is a necessary condition for deformation rigidity''.

\section{First-order deformations and graded rigidity}\label{l09:sec:deform}

The general deformation theory of the library is the subject of Chapter~\ref{ch:L07}, where it is
applied to the Virasoro minimal models. To keep this chapter self-contained we recall every
definition the headline uses, with its Lean source (\leanpath{VOA/Core/Deformation.lean} and
\leanpath{VOA/Core/DeformationGraded.lean}).

\subsection{The idea}

Let $V$ be a vertex operator algebra with state--field map $Y$. Let $R = \C[t]/(t^2)$ be the ring of
dual numbers and $V_R = R\otimes_\C V = V\oplus tV$. A \emph{first-order deformation} of $V$ is a
bilinear map
\begin{equation*}
  Y_1 : V\otimes V \to V((x)),\qquad u\otimes v\mapsto Y_1(u,x)v = \sum_n u^1_{(n)}v\,x^{-n-1},
\end{equation*}
such that $Y^t(u,x) := Y(u,x) + t\,Y_1(u,x)$, extended $R$-bilinearly, makes $V_R$ a vertex algebra
over $R$ with the same vacuum $\lnineone$. Expanding each vertex-algebra axiom in powers of $t$, the
$t^0$ part is the axiom for $V$ and the $t^1$ part is a linear condition on $Y_1$: $Y_1$ is a
\emph{2-cocycle}. A \emph{trivial} deformation is one obtained from $Y$ by an $R$-linear change of
coordinates $1 + t\varphi$ of $V_R$, which gives
\begin{equation}\label{l09:eq:cobdy}
  Y_1(u,x)v = (\delta\varphi)(u,x)v := Y(\varphi u,x)v - \varphi\bigl(Y(u,x)v\bigr) + Y(u,x)\varphi v;
\end{equation}
these are the \emph{coboundaries}. Rigidity is ``every cocycle is a coboundary'', i.e.\ the
vanishing of the second deformation cohomology. Physically: a first-order deformation is an
infinitesimal change of the OPE coefficients compatible with associativity (crossing), and it is
trivial when it is undone by an infinitesimal field redefinition. A rigid chiral algebra has no
infinitesimal moduli at all --- in particular it cannot sit inside a continuous family of chiral
algebras, which is the chiral-algebra analogue of ``no exactly marginal deformations''.

\subsection{The Lean definitions}

\begin{leancode}
structure FirstOrderDeformation (W : VertexOperatorAlgebra 𝒱) where
  Y1 : 𝒱 →ₗ[ℂ] (End ℂ 𝒱)⟦z^±⟧
  truncation : ∀ u v : 𝒱, ∃ N : ℤ, ∀ n : ℤ, N ≤ n → mode (Y1 u) n v = 0
  vacuum_Y1 : Y1 W.vacuum = 0
  creation : ∀ (u : 𝒱) (n : ℤ), -1 ≤ n → mode (Y1 u) n W.vacuum = 0
  jacobi : ∀ (u v x : 𝒱) (ℓ m n : ℤ),
      (∑ᶠ i : ℤ, (negOnePow i * cbz ℓ i) •
          (mode (Y1 u) (ℓ + m - i) (W.smul v (n + i) x)
            + W.smul u (ℓ + m - i) (mode (Y1 v) (n + i) x)
            - negOnePow ℓ •
              (mode (Y1 v) (ℓ + n - i) (W.smul u (m + i) x)
                + W.smul v (ℓ + n - i) (mode (Y1 u) (m + i) x))))
        = ∑ᶠ i : ℤ, cbz m i •
            (W.smul (mode (Y1 u) (ℓ + i) v) (m + n - i) x
              + mode (Y1 (W.smul u (ℓ + i) v)) (m + n - i) x)
\end{leancode}
\begin{leancode}
def smul₁ (D : FirstOrderDeformation W) (u : 𝒱) (n : ℤ) (v : 𝒱) : 𝒱 := mode (D.Y1 u) n v
\end{leancode}

\begin{dictionary}
\textbf{Shape.} \lean{(End ℂ 𝒱)⟦z^±⟧} is the space of two-sided formal series with coefficients in
$\End(V)$ (Chapter~\ref{ch:L01}); \lean{𝒱 →ₗ[ℂ] …} is a $\C$-linear map, so \lean{Y1} is linear in
$u$ and (coefficientwise) in $v$: the bilinear $Y_1$ above. \lean{mode f n = f (-n-1)} (the
coefficient of $z^{-n-1}$), so \lean{D.smul₁ u n v} is $u^1_{(n)}v$.
\textbf{The four fields.} \lean{truncation} is the $t$-part of ``$Y(u,x)v$ has finitely many
negative powers'' (Y1 lands in $V((x))$). \lean{vacuum_Y1} is the $t$-part of $Y^t(\lnineone,x) =
\mathrm{id}$. \lean{creation} is the $t$-part of $u^t_{(n)}\lnineone = 0$ for $n\ge0$ and
$u^t_{(-1)}\lnineone = u$, i.e.\ $u^1_{(n)}\lnineone = 0$ for $n\ge-1$. \lean{jacobi} is the
$t$-linear part of the Borcherds identity (Chapter~\ref{ch:L03}): \lean{cbz ℓ i} is the binomial
coefficient $\binom{\ell}{i}$ ($0$ for $i<0$), \lean{negOnePow i} is $(-1)^i$, \lean{∑ᶠ} is a finite
sum; on each side the single $Y_1$ is inserted in each possible slot (Leibniz rule).
\textbf{Could the sums be vacuous?} Mathlib defines \lean{∑ᶠ} of an infinite-support family to be
$0$, which would make \lean{jacobi} read $0=0$. The library proves both supports finite
(\lean{jacobiLHS_support_finite}, \lean{jacobiRHS_support_finite}), so the axiom is the genuine
identity. \textbf{What is \emph{not} imposed.} The conformal vector, the operator $L_0$ and the
translation operator $T = L_{-1}$ are not required to be preserved; nor is any grading condition.
Gradedness is a separate predicate (next).
\end{dictionary}

\noindent File \leanpath{VOA/Core/Deformation.lean}.
\begin{leancode}
def IsGraded (D : FirstOrderDeformation W) : Prop :=
  ∀ (k ℓ : ℤ) (u : 𝒱), u ∈ W.grading k → ∀ (v : 𝒱), v ∈ W.grading ℓ → ∀ (n : ℤ),
    D.smul₁ u n v ∈ W.grading (k + ℓ - n - 1)
\end{leancode}
\noindent File \leanpath{VOA/Core/Deformation.lean}.
\begin{leancode}
def GradePreserving (W : VertexOperatorAlgebra 𝒱) (φ : 𝒱 →ₗ[ℂ] 𝒱) : Prop :=
  ∀ (k : ℤ) (u : 𝒱), u ∈ W.grading k → φ u ∈ W.grading k
\end{leancode}
\noindent File \leanpath{VOA/Core/Deformation.lean}.
\begin{leancode}
def IsCoboundary (W : VertexOperatorAlgebra 𝒱) (f : 𝒱 →ₗ[ℂ] (End ℂ 𝒱)⟦z^±⟧) : Prop :=
  ∃ φ : 𝒱 →ₗ[ℂ] 𝒱, φ W.vacuum = 0 ∧ GradePreserving W φ ∧
    ∀ (u : 𝒱) (n : ℤ) (v : 𝒱),
      mode (f u) n v = W.smul (φ u) n v - φ (W.smul u n v) + W.smul u n (φ v)
\end{leancode}
\noindent File \leanpath{VOA/Core/Deformation.lean}.
\begin{leancode}
def IsTrivial (D : FirstOrderDeformation W) : Prop := IsCoboundary W D.Y1
\end{leancode}
\noindent File \leanpath{VOA/Core/Deformation.lean}.
\begin{leancode}
def IsRigid (W : VertexOperatorAlgebra 𝒱) : Prop :=
  ∀ D : FirstOrderDeformation W, D.IsTrivial
\end{leancode}
\noindent File \leanpath{VOA/Core/DeformationGraded.lean}.
\begin{leancode}
def IsRigidGraded (W : VertexOperatorAlgebra 𝒱) : Prop :=
  ∀ D : FirstOrderDeformation W, D.IsGraded → D.IsTrivial
\end{leancode}

\begin{definition}\label{l09:def:rigid}
A first-order deformation $D$ is \emph{graded} if $u^1_{(n)}v \in V_{a+b-n-1}$ whenever $u\in V_a$,
$v\in V_b$ (the same weight rule as $u_{(n)}v$). It is \emph{trivial} if $Y_1 = \delta\varphi$ as in
\eqref{l09:eq:cobdy} for some linear $\varphi:V\to V$ that preserves every $V_n$ and kills $\lnineone$.
$V$ is \emph{graded rigid} (in full: \emph{graded first-order rigid}) if every graded first-order
deformation is trivial. In this chapter ``graded rigid'' always means graded first-order rigid.
\end{definition}
\begin{physics}
A graded first-order deformation is an infinitesimal change of all OPE coefficients that keeps conformal
weights fixed and keeps the OPE associative (crossing-symmetric) to first order; a trivial one is an
infinitesimal field redefinition $\phi\mapsto\phi + t\varphi(\phi)$. In a full CFT, deformations are generated by
exactly marginal operators; for a chiral algebra, the analogous question is whether the OPE data have
moduli. In a full CFT with currents $J$, $\bar J$, the current--current operator $J\bar J$ is the standard
exactly marginal deformation (it changes the radius of a free boson); at the purely chiral level, the
simplest non-trivial deformation is moving the level of the universal algebra $V^k$, which is a continuous
parameter there (for a compact WZW model the level is quantised, but the chiral algebra does not know that). Graded rigidity says there are no such moduli.
\end{physics}


\begin{dictionary}
\lean{IsGraded} is literally the weight rule, for all integers $k,\ell,n$. \lean{IsCoboundary} is
\eqref{l09:eq:cobdy} mode by mode: $u^1_{(n)}v = (\varphi u)_{(n)}v - \varphi(u_{(n)}v) +
u_{(n)}(\varphi v)$. The conditions $\varphi\lnineone = 0$ (so the vacuum is unchanged by
$1+t\varphi$) and grade preservation are the standard ones. \lean{IsRigidGraded} is
Definition~\ref{l09:def:rigid}: universally quantified over \emph{all} graded deformations.
\textbf{Comparison with Linshaw--Qi.} In their paper (arXiv:2601.12017v2, \S2) a first-order
deformation is a $Y_1$ making $\C[t]/(t^2)\otimes V$ a grading-restricted vertex algebra \emph{with
the grading operator and the derivative operator preserved}; a deformation is trivial when
$Y_1 = \delta\varphi$ for a grading-preserving $\varphi$ with $\varphi(\lnineone)=0$. The library's
trivial deformations are the same. The library's class of graded deformations is \emph{larger}: it
does not require the translation operator to be preserved ($u^1_{(-2)}\lnineone$ may be nonzero).
Since \lean{IsRigidGraded} asserts triviality for every member of a larger class, \textbf{the Lean
notion of graded rigidity is at least as strong as Linshaw--Qi's} (it implies theirs). It is not
weaker in any respect. Besides $T$, the conformal vector is not part of the data either: a deformation may
make $\omega$ stop being conformal. In standard terms, \lean{IsRigidGraded W} is the vanishing of the graded
second deformation cohomology $H^2$ of $W$ as a vertex algebra with the vacuum fixed --- the same gauge group
as Linshaw--Qi (grade-preserving $\varphi$, $\varphi\lnineone = 0$) and a larger class of deformations. So
\lean{L43_graded_rigidity} below is Linshaw--Qi's Theorem~5.1 in their own sense, not merely an analogue.
\end{dictionary}

\begin{example}[the scaling coboundary]\label{l09:ex:scaling}
For every VOA, $\varphi = L_0$ preserves the grading and kills $\lnineone$, so $\delta(L_0)$ is a
trivial graded deformation. On homogeneous $u\in V_a$, $v\in V_b$,
\begin{equation*}
  \delta(L_0)(u)_{(n)}v = a\,u_{(n)}v - (a+b-n-1)\,u_{(n)}v + b\,u_{(n)}v = (n+1)\,u_{(n)}v .
\end{equation*}
In Lean this is \lean{NonzeroDeform.scalingDeformation} with
\noindent File \leanpath{VOA/Examples/NonzeroDeformation.lean}.
\begin{leancode}
theorem smul₁_scalingDeformation_homog (W : VertexOperatorAlgebra 𝒱) {a b : ℤ} {u v : 𝒱}
    (hu : u ∈ W.grading a) (hv : v ∈ W.grading b) (n : ℤ) :
    (scalingDeformation W).smul₁ u n v = ((n : ℂ) + 1) • W.smul u n v := by
\end{leancode}
This one example will refute two of the hypotheses that an earlier version of the library carried
(Section~\ref{l09:sec:lq}). Note also that $\delta(L_0)$ does not preserve the translation operator
(\lean{smul₁_scaling_conformal_vacuum}: $\omega^1_{(-2)}\lnineone = -L_{-3}\lnineone$), which
illustrates the remark above that the library's deformation class is larger than Linshaw--Qi's.
\end{example}

\begin{warning}
\textbf{\lean{IsRigid} is the wrong target.} \lean{IsRigid} quantifies over \emph{all}
first-order deformations, graded or not, while coboundaries must be grade-preserving. The
coboundary $\delta(L_1)$ of the non-grade-preserving operator $L_1$ is a legitimate (ungraded)
deformation of every VOA, and it is not of the form $\delta\varphi$ with $\varphi$ grade-preserving
whenever $L_1$ acts nontrivially; the library proves this for Virasoro as
\lean{VOA.NonzeroDeform.not_isRigid_virasoroVOA}: \lean{IsRigid} fails for every universal Virasoro
VOA $V(c,0)$, for this definitional reason (Chapter~\ref{ch:L07}). Since a trivial deformation is
automatically graded (\lean{NonzeroDeform.isGraded_of_isTrivial}), any nonzero $\delta(L_1)$ refutes
\lean{IsRigid}; on $L_k(\lninesl)$, $k\neq0$, one computes on paper
$\delta(L_1)(L_{-1}J,-2,\lnineone) = 2L_{-1}J - L_1L_{-1}^2J = -4L_{-1}J\neq0$ for a current $J$, so
\lean{IsRigid} fails there too (this instance is not formalised). \lean{IsRigid} must never be quoted as
true. The headline is therefore stated for \lean{IsRigidGraded}, the graded first-order notion, which is at
least as strong as the notion Linshaw--Qi use.
\end{warning}

\begin{warning}
\textbf{A non-vacuity point the referee should know.} A rigidity theorem is a universal statement,
so it would be worthless if, through some defect, \emph{no} VOA could ever have a non-trivial graded
deformation in this formalism. Three facts speak against that. (i) Deformations exist in abundance:
every coboundary $\delta\varphi$ with $\varphi\lnineone=0$ is one (\lean{NonzeroDeform.deformOf}), so
\lean{FirstOrderDeformation} is not an empty or degenerate type. (ii) The \emph{weaker}
conclusions that would follow from a defect are refuted: the free boson has graded deformations in
``pole form'' with nonzero pole (Section~\ref{l09:sec:lq}, \lean{not_normalisedPoleVanishes_heisenberg}).
(iii) For the universal Virasoro VOA $V(c,0)$, ``$\partial/\partial c$ is a non-trivial graded deformation'' is
proved in full. $Y_1 = \partial_cY$, computed on the $c$-independent PBW basis of the universal $V(c,0)$, is a
graded first-order deformation with $\omega^1_{(1)}\omega = 0$ and $\omega^1_{(3)}\omega = \frac12\lnineone$. In
Lean, every coboundary on a VOA with $V_2 = \C\omega$ has $(\omega^1_{(1)}\omega,\omega^1_{(3)}\omega) =
(2\mu\omega,\mu c\lnineone)$ for one scalar $\mu$ (\lean{VOA.VirCDef.cmode_conformal_of_isCoboundary}), so such a
deformation is non-trivial (\lean{VOA.VirCDef.not_isTrivial_of_omega_witness},
\lean{VOA.VirCDef.not_isTrivial_of_isCDeformation},
\lean{VOA.VirCDef.not_isRigidGraded_of_isCDeformation}; all quoted, with their status, in
Chapter~\ref{ch:L07}, Remark on non-vacuity).
When this chapter was first written, the $c$-deformation was not constructed
(\lean{IsCDeformation c} had no inhabitant), so non-vacuity was argued on paper only. That item is now resolved:
\lean{VOA.}\allowbreak\lean{VirPoly.}\allowbreak\lean{cDeformation} constructs $\partial_cY$ on $V(c,0)$,
\lean{VOA.}\allowbreak\lean{VirPoly.}\allowbreak\lean{isCDeformation_}\allowbreak\lean{cDeformation} proves that it inhabits
\lean{IsCDeformation c}, and \lean{VOA.}\allowbreak\lean{VirPoly.}\allowbreak\lean{not_}\allowbreak\lean{isRigidGraded_}\allowbreak\lean{virasoroVOA} proves
\lean{¬ Deform.IsRigidGraded (virasoroVOA c)} for every $c\in\C$
(\leanpath{VOA/\allowbreak Examples/\allowbreak VirasoroCDefFinal.lean}; Chapter~\ref{ch:L07},
Section~\ref{l07:sec:nonvacuity}). The caveat that remains concerns its scope: $\partial_cY$ does not preserve the central
term of $\omega_{(3)}\omega$ (it moves $c$), which the library's \lean{FirstOrderDeformation} allows; this is the same
convention under which the rigidity theorems of this chapter are proved. (Other non-rigid VOAs certainly exist on paper,
e.g.\ $V^k(\lninesl)$ itself, deformed along $k$; the library does not formalise that one.)
\end{warning}

\subsection{Regularity and the pole}

For weight-one states $a,b$ (currents), a graded deformation has
$a^1_{(n)}b\in V_{1-n}$, which is zero for $n\ge2$ (no negative weights). So the singular part of
$Y_1(a,x)b$ has only two coefficients: $a^1_{(0)}b\in V_1$ (a correction to the Lie bracket) and
$a^1_{(1)}b\in V_0 = \C\lnineone$ (a correction to the level).

\noindent File \leanpath{VOA/Core/Deformation.lean}.
\begin{leancode}
def Regular (D : FirstOrderDeformation W) : Prop :=
  ∀ a ∈ W.grading 1, ∀ b ∈ W.grading 1, ∀ n : ℤ, 0 ≤ n → D.smul₁ a n b = 0
\end{leancode}
\noindent File \leanpath{VOA/Core/Deformation.lean}.
\begin{leancode}
def poleVec (D : FirstOrderDeformation W) (a b : 𝒱) : 𝒱 := D.smul₁ a 1 b
\end{leancode}
\begin{leancode}
def FirstOrderDeformation.InPoleForm (D : FirstOrderDeformation W) : Prop :=
  ∀ a ∈ W.grading 1, ∀ b ∈ W.grading 1, D.smul₁ a 0 b = 0
\end{leancode}
\noindent File \leanpath{VOA/Examples/AffineSl2Rigidity43.lean}.
\begin{leancode}
theorem regular_iff_zero_one {𝒱 : Type*} [AddCommGroup 𝒱] [Module ℂ 𝒱]
    {W : VertexOperatorAlgebra 𝒱} (D : Deform.FirstOrderDeformation W) (hD : D.IsGraded) :
    D.Regular ↔ ∀ a ∈ W.grading 1, ∀ b ∈ W.grading 1,
      D.smul₁ a 0 b = 0 ∧ D.smul₁ a 1 b = 0 := by
\end{leancode}

\begin{dictionary}
\lean{Regular} is ``$Y_1(a,x)b$ has no negative powers of $x$ for weight-one $a,b$'': in the
index convention \lean{mode f n = f(-n-1)}, negative powers $x^{-n-1}$ are exactly $n\ge0$ (the
library also proves the literal coefficient form, \lean{regular_iff}). Since $a = a_{(-1)}\lnineone$,
$Y_1(a,x)b$ is Linshaw--Qi's $Y_1(a(-1)\lnineone,x)b(-1)\lnineone$, and quantifying over all of $V_1$
is equivalent to quantifying over their generators by bilinearity. \lean{poleVec a b} is the
$x^{-2}$ coefficient $a^1_{(1)}b$ (proved \lean{poleVec_eq_coeff}); for graded $D$ it lies in
$V_0$ (\lean{poleVec_mem_grading_zero}). \lean{InPoleForm} says the $x^{-1}$ coefficient vanishes on
currents: the first clause of Linshaw--Qi's equation (3). \lean{regular_iff_zero_one} is the
mode count just explained (VOA positivity kills $n\ge2$).
\end{dictionary}

\section{Weight-one cohomology: normalising a deformation}\label{l09:sec:whitehead}

\leanpath{VOA/Core/WeightOneCohomology.lean} proves the step that Linshaw--Qi perform silently in
their equation (3) (citing Kovalchuk--Qi \S6.2): every graded deformation can be changed by a
coboundary so that the $x^{-1}$ coefficient on currents vanishes. It is pure Lie theory of $V_1$.

\subsection{Chevalley--Eilenberg cochains in degree two}

For a Lie algebra $\lnineg$ acting on itself by the adjoint action, a \emph{1-cochain} is a linear
map $\psi:\lnineg\to\lnineg$ and a \emph{2-cochain} is an alternating bilinear map
$\omega:\lnineg\times\lnineg\to\lnineg$. The differentials are
\begin{align*}
  (d_{12}\psi)(x,y) &= [x,\psi y] - [y,\psi x] - \psi[x,y],\\
  (d_{23}\omega)(x,y,z) &= [x,\omega(y,z)] - [y,\omega(x,z)] + [z,\omega(x,y)]
     - \omega([x,y],z) + \omega([x,z],y) - \omega([y,z],x).
\end{align*}
$\omega$ is a \emph{cocycle} if $d_{23}\omega = 0$; $H^2(\lnineg,\lnineg) = $ cocycles$/d_{12}($1-cochains$)$.
Mathlib provides the cochain spaces and these two differentials
(\lean{LieModule.Cohomology.twoCochain}, \lean{d₁₂}, \lean{d₂₃}, \lean{twoCocycle}, by Carnahan) but no
cohomology groups and no Whitehead lemma. The library therefore states vanishing as ``every
cocycle is a coboundary'':

\begin{leancode}
def WeightOneH2Vanishes (W : VertexOperatorAlgebra 𝒱) : Prop :=
  ∀ ω : twoCochain ℂ ↥(W.grading 1) ↥(W.grading 1),
    ω ∈ twoCocycle ℂ ↥(W.grading 1) ↥(W.grading 1) →
      ∃ ψ : oneCochain ℂ ↥(W.grading 1) ↥(W.grading 1),
        d₁₂ ℂ ↥(W.grading 1) ↥(W.grading 1) ψ = ω
\end{leancode}
\begin{leancode}
def CentrelessWeightOne (W : VertexOperatorAlgebra 𝒱) : Prop :=
  ∀ z : ↥(W.grading 1), (∀ x : ↥(W.grading 1), ⁅z, x⁆ = 0) → z = 0
\end{leancode}

\begin{dictionary}
The Lie algebra is $V_1$ with bracket \lean{⁅u, v⁆} $= u_{(0)}v$ (\lean{coe_bracket_gradingOne};
the Lie structure on $V_1$ is Chapter~\ref{ch:L04}). \lean{twoCochain ℂ L M} is Mathlib's space of
alternating bilinear maps $L\times L\to M$, \lean{oneCochain} the linear maps, and \lean{d₁₂},
\lean{twoCocycle} as displayed above (the file docstring checks that Mathlib's sign convention is
$d_{12}\psi(x,y) = [x,\psi y]-[y,\psi x]-\psi[x,y]$). So \lean{WeightOneH2Vanishes W} is
$H^2(V_1,V_1) = 0$ and \lean{CentrelessWeightOne W} is $Z(V_1) = 0$. Neither mentions any
deformation, any mode other than $(0)$, or anything outside $V_1$: they are invariants of the Lie
algebra $V_1$ alone.
\end{dictionary}

\begin{theorem}[Whitehead's second lemma for $\lninesl$]\label{l09:thm:whitehead}
If a complex Lie algebra $\lnineg$ has a basis $(e,h,f)$ forming an $\lninesl$-triple, then every
2-cocycle $\omega\in Z^2(\lnineg,\lnineg)$ is $d_{12}\psi$ for some $\psi$; and $Z(\lnineg)=0$.
\end{theorem}
\begin{physics}
$H^2(\lnineg,\lnineg)=0$ is the statement that a semisimple Lie algebra has no infinitesimal deformations of
its bracket: the symmetry algebra of the currents cannot be continuously deformed. Physically, the
single-pole term $J^{[a,b]}/(z-w)$ of the current OPE cannot be changed except by a rotation of the current
basis.
\end{physics}


\begin{leancode}
theorem whitehead2_of_sl2Basis (b : Basis (Fin 3) ℂ g) (ht : IsSl2Triple (b 1) (b 0) (b 2))
    (ω : twoCochain ℂ g g) (hω : ω ∈ twoCocycle ℂ g g) :
    ∃ ψ : oneCochain ℂ g g, d₁₂ ℂ g g ψ = ω := by
\end{leancode}
\begin{leancode}
theorem weightOneH2Vanishes_of_sl2Basis (W : VertexOperatorAlgebra 𝒱)
    (b : Basis (Fin 3) ℂ ↥(W.grading 1)) (ht : IsSl2Triple (b 1) (b 0) (b 2)) :
    WeightOneH2Vanishes W :=
  fun ω hω => Sl2Whitehead.whitehead2_of_sl2Basis b ht ω hω
\end{leancode}

\begin{proof}[Proof on paper]
Write $\omega(e,h) = A_1e+A_2h+A_3f$, $\omega(e,f) = B_1e+B_2h+B_3f$, $\omega(h,f) = C_1e+C_2h+C_3f$.
Evaluate $d_{23}\omega(e,h,f) = 0$. The three terms $\omega([\cdot,\cdot],\cdot)$ cancel identically:
$[e,h] = -2e$ and $[h,f] = -2f$ contribute $+2\omega(e,f)$ and $-2\omega(e,f)$, and $[e,f] = h$
contributes $\omega(h,h)=0$. What remains, $[e,\omega(h,f)] - [h,\omega(e,f)] + [f,\omega(e,h)] = 0$,
expands in the basis to exactly three scalar equations,
\begin{equation*}
  C_2 + B_1 = 0,\qquad C_3 = A_1,\qquad B_3 + A_2 = 0 .
\end{equation*}
Now take
\begin{equation*}
  \psi e = \bigl(B_2 - \tfrac{A_1}{2}\bigr)e + \tfrac{A_3}{4}f,\qquad
  \psi h = C_2 e - \tfrac{A_1}{2}h + A_2 f,\qquad
  \psi f = \tfrac{C_1}{4}e .
\end{equation*}
For general $\psi e = p_1e+p_2h+p_3f$, $\psi h = q_1e+q_2h+q_3f$, $\psi f = r_1e+r_2h+r_3f$ one computes
\begin{align*}
  d_{12}\psi(e,h) &= -2q_2\,e + (q_3+2p_2)\,h + 4p_3\,f,\\
  d_{12}\psi(e,f) &= -(2r_2+q_1)\,e + (p_1+r_3-q_2)\,h - (2p_2+q_3)\,f,\\
  d_{12}\psi(h,f) &= 4r_1\,e + (q_1+2r_2)\,h - 2q_2\,f,
\end{align*}
and substituting the chosen $\psi$ and the three cocycle relations reproduces $\omega$ on the three
basis pairs, hence everywhere. The solution is unique up to a three-parameter family, which is
$\lnineDer(\lninesl) = \lninead(\lninesl)$. Centrelessness: if $[z,x] = 0$ for all $x$, then $[z,f] = 0$
kills the $e$- and $h$-coefficients of $z$, and $[z,e]=0$ kills the $f$-coefficient.
\end{proof}

\subsection{The bridge: a graded deformation gives a 2-cocycle}

For a graded deformation, $a^1_{(0)}b\in V_{1+1-0-1} = V_1$ for $a,b\in V_1$, so
$\omega(a,b) := a^1_{(0)}b$ is a bilinear map $V_1\times V_1\to V_1$. The $\ell=m=n=0$ case of the
deformed Jacobi identity reads
\begin{equation}\label{l09:eq:jacobi000}
  \omega(u,[v,x]) + [u,\omega(v,x)] - \omega(v,[u,x]) - [v,\omega(u,x)]
  = [\omega(u,v),x] + \omega([u,v],x).
\end{equation}
Setting $u=v$ gives $[\omega(a,a),x] = 0$: $\omega(a,a)$ is central, hence zero when $V_1$ is
centreless, so $\omega$ is alternating. Using that, \eqref{l09:eq:jacobi000} rearranges to
$d_{23}\omega = 0$.

\noindent File \leanpath{VOA/Core/WeightOneCohomology.lean}.
\begin{leancode}
theorem weightOneCochain_mem_twoCocycle (D : FirstOrderDeformation W) (hD : D.IsGraded)
    (hZ : W.CentrelessWeightOne) :
    weightOneCochain D hD hZ ∈ twoCocycle ℂ ↥(W.grading 1) ↥(W.grading 1) := by
\end{leancode}
\noindent File \leanpath{VOA/Core/WeightOneCohomology.lean}.
\begin{leancode}
theorem coboundary_weightOne_eq_d₁₂ (ψ : ↥(W.grading 1) →ₗ[ℂ] ↥(W.grading 1))
    (a b : ↥(W.grading 1)) :
    W.smul (extendWeightOne W ψ (a : 𝒱)) 0 (b : 𝒱)
        - extendWeightOne W ψ (W.smul (a : 𝒱) 0 (b : 𝒱))
        + W.smul (a : 𝒱) 0 (extendWeightOne W ψ (b : 𝒱))
      = ((d₁₂ ℂ ↥(W.grading 1) ↥(W.grading 1) ψ a b : ↥(W.grading 1)) : 𝒱) := by
\end{leancode}

\begin{dictionary}
\lean{weightOneCochain D hD hZ} is $\omega(a,b) = a^1_{(0)}b$ packaged as an alternating cochain;
its alternation proof uses the centrelessness hypothesis \lean{hZ}, exactly as in the paper argument
(the file warns that skew-symmetry of $\omega$ is \emph{not} free --- the alternative would be a port
of the whole skew-symmetry theory to the deformed setting). \lean{extendWeightOne W ψ} is the
extension of $\psi:V_1\to V_1$ by zero on the other weight spaces, which is grade-preserving and
kills the vacuum (\lean{gradePreserving_extendWeightOne}, \lean{extendWeightOne_vacuum}). The second
theorem says that the VOA coboundary $\delta\varphi$ of that extension, restricted to mode $0$ on
$V_1$, is literally Mathlib's $d_{12}\psi$: the signs agree on the nose.
\end{dictionary}

\begin{corollary}[normalisation]\label{l09:cor:normalise}
If $V_1\cong\lninesl$, every graded deformation $D$ differs by a graded coboundary from a graded
deformation $E$ in pole form ($a^1_{(0)}b = 0$ on $V_1$).
\end{corollary}
\begin{physics}
This is a gauge fixing: by redefining the currents one can always assume the deformation does not touch
the single pole of $J^aJ^b$, i.e.\ keeps the symmetry algebra $SU(2)$ fixed. What is left is the double
pole: a possible change of the level $k\mapsto k+\delta t$, the only physically meaningful parameter.
\end{physics}


\begin{proof}[Proof on paper]
By the bridge, $\omega = a^1_{(0)}b$ is a cocycle; by Theorem~\ref{l09:thm:whitehead}, $\omega =
d_{12}\psi$; with $\varphi$ the extension of $\psi$ by zero, $E := D - \delta\varphi$ is graded (coboundaries
of grade-preserving maps are graded) and its mode-$0$ part on $V_1$ is $\omega - d_{12}\psi = 0$.
\end{proof}

\paragraph{Why the normalisation matters.} Before normalisation, the pole coefficient
$a^1_{(1)}b$ is not an invariant of the deformation class: the trivial deformation $\delta(L_0)$ of
Example~\ref{l09:ex:scaling} has $a^1_{(1)}b = 2a_{(1)}b = 2k\kappa(a,b)\lnineone$, a nonzero ``level
change'' that is pure gauge. After normalisation the residual gauge freedom preserving
$a^1_{(0)}b = 0$ consists of $\varphi$ with $\varphi|_{V_1}\in\lnineDer(V_1) = \lninead(V_1)$, and invariance of
$\kappa$ shows these do not change the pole. So only after normalisation is ``the level change
$\delta$'' a well-defined quantity. This is the conceptual heart of the split of Linshaw--Qi's
theorem in the next section.

\section{Linshaw--Qi Theorem 5.1 at \texorpdfstring{$k=-4/3$}{k=-4/3}, and the bundled hypothesis}\label{l09:sec:lq}

\subsection{The source}

A.~R.~Linshaw and F.~Qi, \emph{Deformation rigidity of some simple affine VOAs},
arXiv:2601.12017 (v2), prove two rigidity theorems: Theorem~4.1, every simple affine VOA at a
positive integer level admits only trivial deformations; and Theorem~5.1, the simple affine VOA
$L_{-4/3}(\lninesl)$ admits only trivial deformations. Both rest on their Theorem~3.1:
\emph{if $V$ is a simple affine VOA and $Y_1(a(-1)\lnineone,x)b(-1)\lnineone$ is regular (has no negative
powers of $x$) for all $a,b\in\lnineg$, then $Y_1$ is trivial.} The proof of Theorem~5.1 then shows,
using the weight-three singular vector at $k=-4/3$, that the singular part of $Y_1$ on currents
vanishes. They conjecture that rigidity holds for every simple affine $L_k(\lnineg)$ that differs
from $V^k(\lnineg)$.

The library contains two developments of Theorem~5.1: a \emph{split} conditional assembly
(file \leanpath{AffineSl2Rigidity43.lean}), which was the vehicle for discovering and
repairing two defective formalisations, and the final \emph{unconditional} proof
(file \leanpath{AffineSl2RigidityComplete.lean} and companions), both in \leanpath{VOA/Examples/}. Both are documented here,
because the first is where the ``tautological bundled hypothesis'' mentioned in the library map
lives.

\subsection{The object \texorpdfstring{$L_{-4/3}(\mathfrak{sl}_2)$}{L(-4/3)}}

\begin{leancode}
noncomputable def L43 :
    VertexOperatorAlgebra (WakModule (-4 / 3 : ℂ) ⧸ QuotVOA.maxGradedIdeal voaSl2Adm) :=
  QuotVOA.maxQuotientVOA voaSl2Adm
\end{leancode}
\noindent File \leanpath{VOA/Examples/AffineSl2VOAInstance.lean}.
\begin{leancode}
noncomputable def voaSl2Adm : VertexOperatorAlgebra (WakModule (-4 / 3 : ℂ)) :=
  affineSl2VOA_closed (-4 / 3 : ℂ) admLevel_noncritical
\end{leancode}

\begin{dictionary}
\lean{voaSl2Adm} is the universal $V^{-4/3}(\lninesl)$ (\lean{admLevel_noncritical} is the proof of
$-4/3+2\neq0$), and \lean{L43} its quotient by the maximal proper graded ideal, i.e.\ exactly
Definition~\ref{l09:def:Lk} at $k=-4/3$. It is definitionally the same object as
\lean{AffineSl2GeneralRigidity.simple (-4/3) _} (the two differ only in the proof of
non-criticality, and proofs are irrelevant), so the uniform theorem of Section~\ref{l09:sec:headline}
at $(p,q) = (2,3)$ is about the same algebra --- up to the propositional, not syntactic, equality
$2/3 - 2 = -4/3$ of the level.
\end{dictionary}

Non-vacuity of anything said about \lean{L43} requires that it be nonzero, simple, and a
\emph{proper} quotient of $V^{-4/3}$ (otherwise one would be proving rigidity of the universal
algebra, which is false). All three are proved, $\lean{pbw}$-free:

\begin{leancode}
theorem L43_isSimple : L43.IsSimple :=
  QuotVOA.isSimple_maxQuotientVOA' voaSl2Adm vacuum_ne_zero
\end{leancode}
\begin{leancode}
theorem maxGradedIdeal_ne_bot : QuotVOA.maxGradedIdeal voaSl2Adm ≠ ⊥ := by
\end{leancode}
\begin{leancode}
theorem mkQ_sl2Singular_eq_zero :
    (QuotVOA.maxGradedIdeal voaSl2Adm).mkQ Sl2Sing.sl2Singular = 0 := by
\end{leancode}
The singular vector is the one printed by Linshaw--Qi (after Adamovi\'c),
\begin{equation}\label{l09:eq:sing43}
  s = -48\,h_{-1}e_{-2}\lnineone + 36\,e_{-1}^2f_{-1}\lnineone - 6\,h_{-2}e_{-1}\lnineone
      + 9\,h_{-1}^2e_{-1}\lnineone + 80\,e_{-3}\lnineone ,
\end{equation}
and the library proves it equal to its own independently derived vector:
\begin{leancode}
theorem lqVector_eq_sl2Singular : lqVector (-4 / 3 : ℂ) = Sl2Sing.sl2Singular := by
\end{leancode}

\begin{proof}[Proof on paper of properness]
The canonical invariant bilinear form $(\cdot|\cdot)$ on $V^k$ (normalised by $(\lnineone|\lnineone)=1$)
exists because $L_1V_1 = 0$ (the currents are quasi-primary). For a weight-one quasi-primary $u$,
invariance collapses to $(u_{(n)}a\,|\,b) = -(a\,|\,u_{(-n)}b)$. Every weight-3 vector is a combination of
states $J^a_{-m}y$ with $m\ge1$; hence $(J^a_{-m}y\,|\,s) = -(y\,|\,J^a_ms) = 0$ because $s$ is
annihilated by all positive modes. States of other weights are orthogonal to $s$ for weight reasons.
So $s$ lies in the radical of the form, which by Li's theorem is the maximal graded ideal. Finally
$s\neq0$, proved in the Wakimoto free-field realisation (\lean{Wak.sl2Singular_ne_zero}).
\end{proof}

\begin{warning}
\textbf{A documentation error corrected.} An older file (\leanpath{Examples/AffineSl2Singular.lean} \S D)
claimed a ``discrepancy with the reported literature form'' of \eqref{l09:eq:sing43}. That was a
miscount: it reordered $h_{-1}e_{-2}\lnineone = e_{-2}h_{-1}\lnineone + 2e_{-3}\lnineone$ but forgot that
$9h_{-1}^2e_{-1}\lnineone$ also contributes to the $e_{-2}h_{-1}\lnineone$ monomial
($h_{-1}^2e_{-1}\lnineone = M_5+4M_2+4M_1$). \lean{lqVector_eq_sl2Singular} is the machine refutation;
there is no discrepancy.
\end{warning}

\subsection{The linear algebra of Linshaw--Qi \S5.3, and an erratum}

Linshaw--Qi reduce Theorem~5.1 to five linear equations (10)--(14) in sixteen unknowns: fifteen
coordinates $a_i,b_i,c_i$ of three unknown weight-3 vectors and the pole constant $c$. The library
transcribes them twice:

\noindent File \leanpath{VOA/Examples/AffineSl2Rigidity43.lean}.
\begin{leancode}
structure CorrectedSystem (a₁ a₂ a₃ a₄ a₅ b₁ b₂ b₃ b₄ b₅ c₁ c₂ c₃ c₄ c₅ c : ℂ) : Prop where
  eq10 : 84 * a₃ + 168 * a₄ - 28 * a₅ - 18 * b₃ - 36 * b₄ + 6 * b₅
      - 72 * c₃ - 144 * c₄ + 24 * c₅ + 72 * c = 0
  eq11 : -42 * a₂ - 56 * a₄ + 9 * b₂ + 12 * b₄ + 36 * c₂ + 48 * c₄ + 9 * c = 0
  eq12 : -42 * a₁ - 56 * a₃ + 9 * b₁ + 12 * b₃ + 36 * c₁ + 48 * c₃ - 6 * c = 0
  eq13 : 84 * a₂ - 560 * a₄ + 168 * a₅ - 18 * b₂ + 120 * b₄ - 36 * b₅
      - 72 * c₂ + 480 * c₄ - 144 * c₅ + 36 * c = 0
  eq14 : 84 * a₁ - 280 * a₂ - 168 * a₃ + 784 * a₄ - 336 * a₅ - 18 * b₁ + 60 * b₂ + 36 * b₃
      - 168 * b₄ + 72 * b₅ - 72 * c₁ + 240 * c₂ + 144 * c₃ - 672 * c₄ + 288 * c₅ - 96 * c = 0
\end{leancode}
\noindent File \leanpath{VOA/Examples/AffineSl2Rigidity43.lean}.
\begin{leancode}
theorem c_eq_zero_certificate : (1890 : ℂ) * c = 0 := by
  linear_combination 30 * h.eq10 - 14 * h.eq11 + 18 * h.eq12 + 23 * h.eq13 + 9 * h.eq14
\end{leancode}
(and \lean{PrintedSystem}, identical except that \lean{eq10} ends in \lean{12 * c}).

\begin{warning}
\textbf{Erratum in the source.} Equation (10) is printed with constant term $+12c$; the paper's own
\S5.1 display gives $+72c$, and the fifteen other coefficients of (10) are correct. The printed
elimination is internally consistent with the wrong $12$ and ends with $10c = 0$; the corrected one
gives $210c=0$. The conclusion $c=0$ is unaffected. The one-line certificate: the covector
$\lambda = (30,-14,18,23,9)$ annihilates all fifteen $a_i,b_i,c_i$ columns, so $\lambda\cdot(10,\ldots,14)$
is $1890c = 0$ (corrected) or $90c=0$ (printed). Moreover the unknowns enter only through
$P_i = 42a_i - 9b_i - 36c_i$, and the solution space is $c=0$ with $(P_1,\ldots,P_5)\in\C\cdot(8,-12,-6,9,36)$
--- the coordinates of $s$ itself in the paper's basis (\lean{PrintedSystem.reduced},
\lean{sl2Singular_eq_lqBasis}). This is an internal consistency check the paper does not make. A
further caution recorded in the file: (10)--(12) are components in $\{e_{-1}f_{-1}\lnineone, h_{-2}\lnineone,
h_{-1}h_{-1}\lnineone\}$ but (13)--(14) in $\{h_{-1}e_{-1}\lnineone, e_{-2}\lnineone\}$, a silent switch of basis.
\end{warning}

These theorems are statements about complex numbers; they say nothing about deformations until
the five equations are derived from a deformation, which is the hard part.

\subsection{The split assembly}

The library's conditional form of Theorem~5.1 separates three inputs:

\noindent File \leanpath{VOA/Examples/AffineSl2Rigidity43.lean}.
\begin{leancode}
def RegularImpliesTrivialGraded {𝒱 : Type*} [AddCommGroup 𝒱] [Module ℂ 𝒱]
    (W : VertexOperatorAlgebra 𝒱) : Prop :=
  ∀ D : Deform.FirstOrderDeformation W, D.IsGraded → D.Regular → D.IsTrivial
\end{leancode}
\noindent File \leanpath{VOA/Examples/AffineSl2Rigidity43.lean}.
\begin{leancode}
def NormalisedPoleVanishes {𝒱 : Type*} [AddCommGroup 𝒱] [Module ℂ 𝒱]
    (W : VertexOperatorAlgebra 𝒱) : Prop :=
  ∀ D : Deform.FirstOrderDeformation W, D.IsGraded → D.InPoleForm →
    ∀ a ∈ W.grading 1, ∀ b ∈ W.grading 1, D.poleVec a b = 0
\end{leancode}
\noindent File \leanpath{VOA/Examples/AffineSl2Rigidity43.lean}.
\begin{leancode}
theorem isRigidGraded_of_split {𝒱 : Type*} [AddCommGroup 𝒱] [Module ℂ 𝒱]
    (W : VertexOperatorAlgebra 𝒱) (hZ : W.CentrelessWeightOne)
    (h31 : RegularImpliesTrivialGraded W) (h1 : W.WeightOneH2Vanishes)
    (h2 : NormalisedPoleVanishes W) : Deform.IsRigidGraded W := by
\end{leancode}
\noindent File \leanpath{VOA/Examples/AffineSl2Rigidity43.lean}.
\begin{leancode}
theorem isRigidGraded_of_sl2Basis {𝒱 : Type*} [AddCommGroup 𝒱] [Module ℂ 𝒱]
    (W : VertexOperatorAlgebra 𝒱) (hb : Sl2BasisOnWeightOne W)
    (h31 : RegularImpliesTrivialGraded W) (h2 : NormalisedPoleVanishes W) :
    Deform.IsRigidGraded W :=
  isRigidGraded_of_split W (centrelessWeightOne_of_sl2Basis' W hb) h31
    (weightOneH2Vanishes_of_sl2Basis' W hb) h2
\end{leancode}
\noindent File \leanpath{VOA/Examples/AffineSl2Rigidity43.lean}.
\begin{leancode}
theorem L43_isRigidGraded (hb : Sl2BasisOnWeightOne L43) (h31 : RegularImpliesTrivialGraded L43)
    (h2 : NormalisedPoleVanishes L43) : Deform.IsRigidGraded L43 :=
  isRigidGraded_of_sl2Basis L43 hb h31 h2
\end{leancode}

\begin{theorem}[split form of Linshaw--Qi Theorem 5.1]\label{l09:thm:split}
For any VOA $W$ with centreless $V_1$: if (Theorem~3.1) every graded regular deformation is trivial,
(Piece~1) $H^2(V_1,V_1) = 0$, and (Piece~2) every graded deformation in pole form has vanishing pole
on $V_1$, then $W$ is graded rigid.
\end{theorem}
\begin{physics}
In physics language: (Piece~1) the symmetry algebra cannot change; (Piece~2) the level cannot change; and
(Theorem~3.1) once both current OPEs are unchanged, the whole chiral algebra is unchanged, because it is
generated by the currents (states are normal-ordered products of currents, and their OPEs follow from the
current OPE by Wick-type rules). The split separates the Lie-theoretic and the level-dependent input.
\end{physics}


\begin{proof}[Proof on paper]
Given graded $D$, Corollary~\ref{l09:cor:normalise} gives $E = D-\delta\varphi$ graded in pole form; Piece~2
kills its pole, so $E$ is regular by \lean{regular_iff_zero_one}; Theorem~3.1 makes $E$ trivial; and
$D = E + \delta\varphi$ is a sum of coboundaries.
\end{proof}

\begin{dictionary}
\lean{RegularImpliesTrivialGraded W} is Linshaw--Qi's Theorem~3.1 \emph{as a property of a
specific} $W$ (with the gradedness that the paper assumes throughout made explicit, since in Lean
gradedness is a predicate). \lean{NormalisedPoleVanishes W} is their \S5 computation: its input is a
deformation \emph{already normalised} (\lean{InPoleForm}), its conclusion is that the pole
\lean{poleVec} vanishes. The three hypotheses of \lean{L43_isRigidGraded} are all now
\emph{theorems} at \lean{L43}: \lean{L43_sl2BasisOnWeightOne}, \lean{L43_regularImpliesTrivialGraded},
and \lean{L43_normalisedPoleVanishes} (next subsection), so the conditional form has been fully
discharged.
\end{dictionary}

\subsection{The two dead ends: a vacuous hypothesis and a tautological one}\label{l09:sec:deadends}

The split exists because two earlier formalisations of the same theorem were wrong in instructive
ways. Both are recorded in \leanpath{VOA/Examples/AffineSl2RigidityCollapse.lean} as theorems.

\begin{warning}
\textbf{Dead end 1: a refutable hypothesis (vacuous headline).} The first version took as its \S5
hypothesis
\begin{leancode}
def GradedImpliesRegular (W : VertexOperatorAlgebra 𝒱) : Prop :=
  ∀ D : Deform.FirstOrderDeformation W, D.IsGraded → D.Regular
\end{leancode}
``every graded deformation is regular''. This is \textbf{false} at $L_{-4/3}$, so a theorem
assuming it proves nothing. The counterexample is the trivial deformation $\delta(L_0)$ of
Example~\ref{l09:ex:scaling}: at $n=0$ it has $a^1_{(0)}b = a_{(0)}b$, and $h_{(0)}e = 2e\ne0$. Even
the weakening \lean{GradedImpliesPoleZero} (``the pole vanishes for every graded deformation'')
fails: at $n=1$, $f^1_{(1)}e = 2f_{(1)}e = 2k\lnineone\neq 0$.
\begin{leancode}
theorem not_gradedImpliesRegular_L43 : ¬ GradedImpliesRegular LQ43.L43 := by
\end{leancode}
\begin{leancode}
theorem not_gradedImpliesPoleZero_L43 : ¬ GradedImpliesPoleZero LQ43.L43 := by
\end{leancode}
The lesson: coboundaries are graded cocycles and are generically not regular, so no statement of
the form ``every graded cocycle is regular'' can hold. Linshaw--Qi do not claim it --- their
equation (3) already presupposes the normalisation of Section~\ref{l09:sec:whitehead}.
\end{warning}

\begin{warning}
\textbf{Dead end 2: the tautological bundled hypothesis.} The repair was
the predicate \lean{GradedRegularUpToCoboundary}, written $\mathrm{GRUC}(W)$ below:
\begin{equation*}
  \mathrm{GRUC}(W):\quad \forall D \text{ graded},\ \exists E \text{ graded, regular, with }
  Y_1^D - Y_1^E \text{ a coboundary}.
\end{equation*}
This is true and non-vacuous, but it is \textbf{equivalent to the conclusion for a trivial reason}:
if $W$ is graded rigid, take $E = 0$. The derivation $\mathrm{IsRigidGraded}(W)\Rightarrow
\mathrm{GRUC}(W)$ uses nothing about $W$, and the file at the time literally
proved the equivalence. A headline ``$\mathrm{GRUC}(L_{-4/3}) \Rightarrow
\mathrm{IsRigidGraded}(L_{-4/3})$'' is bookkeeping: it bundles the normalisation (Piece~1, pure Lie
theory of $V_1$) with the computation (Piece~2, where the level and the singular vector enter) into
one existential statement whose content is the theorem itself. This is what the library map calls
the ``tautological bundled hypothesis''. The predicate and all its supporting lemmas
(\lean{GradedRegularUpToCoboundary}, \lean{trivial_satisfies}, \lean{gradedRegularUpToCoboundary_iff_isRigidGraded}, \ldots)
were \emph{deleted}, not merely superseded; they did not exist at the documented commit, and a pre-release build commit has
no such declarations either (the names survive only in a docstring of \leanpath{AffineSl2Rigidity43.lean} that
records this dead end).
\end{warning}

\paragraph{What ``non-tautological'' can mean.} At a fixed $W$, \emph{any} true hypothesis is
logically equivalent to any true conclusion, so ``not implied by the conclusion'' is not a usable
criterion. The criterion with content is \emph{uniform derivability}: is there a proof of
$\forall W,\ \mathrm{IsRigidGraded}(W)\Rightarrow H(W)$ that uses nothing about $W$? Because rigidity
turns every graded $D$ into a coboundary, this amounts to asking whether the \emph{coboundary
content} of $H$ (its restriction to $D = \delta\varphi$) holds for every $W$. For the deleted predicate
it did (take $E=0$). For the split pieces it does not, and the library proves this:

\begin{leancode}
theorem not_normalisedPoleVanishes_heisenberg (k : ℂ) (hk : k ≠ 0) :
    ¬ LQ43.NormalisedPoleVanishes (heisenbergVOA k hk) := fun h =>
\end{leancode}
\begin{leancode}
theorem not_trivial_satisfies_piece2 :
    ¬ (∀ (𝒲 : Type) [AddCommGroup 𝒲] [Module ℂ 𝒲] (V : VertexOperatorAlgebra 𝒲)
        (D : Deform.FirstOrderDeformation V), D.IsTrivial → D.IsGraded → D.InPoleForm →
        ∀ a ∈ V.grading 1, ∀ b ∈ V.grading 1, D.poleVec a b = 0) := fun H =>
\end{leancode}
\begin{leancode}
theorem no_uniform_derivation_piece2 (k : ℂ) (hk : k ≠ 0)
    (H : ∀ (𝒲 : Type) [AddCommGroup 𝒲] [Module ℂ 𝒲] (V : VertexOperatorAlgebra 𝒲),
      Deform.IsRigidGraded V → LQ43.NormalisedPoleVanishes V) :
    ¬ Deform.IsRigidGraded (heisenbergVOA k hk) :=
\end{leancode}

\begin{dictionary}
At the rank-one free boson (Heisenberg VOA at level $k\neq0$, Chapter~\ref{ch:L05}), $V_1 = \C\alpha$ is
abelian, so $\delta(L_0)$ is automatically in pole form ($a^1_{(0)}b = a_{(0)}b = 0$), yet its pole is
$\alpha^1_{(1)}\alpha = 2k\lnineone\ne0$. Hence Piece~2 is \emph{false} for the free boson, and its
coboundary content (\lean{not_trivial_satisfies_piece2}, the exact analogue of the deleted
\lean{trivial_satisfies}) is refutable. Consequently a uniform derivation of Piece~2 from rigidity
would refute graded rigidity of the free boson (\lean{no_uniform_derivation_piece2}). Physically:
for a free boson a level change is pure gauge (rescale $\alpha$), whereas for $\lninesl$ the bracket
$[h,e]=2e$ pins the normalisation of the currents, so the level cannot be rescaled away. Piece~2
thus encodes a genuine, $W$-specific fact. The same $\delta(L_0)$ cannot refute Piece~2 at
$L_{-4/3}$, because there it is not in pole form ($h_{(0)}e\neq0$):
\lean{scalingDeformation_not_inPoleForm_L43}. Similarly \lean{no_uniform_derivation_sl2Basis}
(instantiated at the triplet algebra $\mathcal{W}(2)$, whose $V_1=0$) shows the $\lninesl$-basis
hypothesis is not uniformly derivable from rigidity either.
\end{dictionary}

\begin{warning}
\textbf{The residual, stated plainly (from the source).} At $L_{-4/3}$ and \emph{given}
$V_1\cong\lninesl$, Piece~2 is still mathematically \emph{equivalent} to graded rigidity: rigidity turns
a deformation in pole form into $\delta\varphi$ with $\varphi|_{V_1}$ a derivation; Whitehead's first lemma
makes it $\lninead x$; invariance of $\kappa$ kills the pole. This is unavoidable --- the conclusion of
\S5 \emph{is} ``the normalised pole vanishes'' --- and it is not the repaired defect: this
derivation consumes the $\lninesl$ structure and a second Whitehead lemma, rather than using nothing
about $W$. In any case the point is now moot for the headline results: the conditional split is no
longer needed, because \lean{L43_graded_rigidity} below and the uniform theorem prove rigidity
\emph{without} assuming Piece~2, by deriving \S5's \emph{input} (the singular-vector relation) from
the definition of $L_k$.
\end{warning}

\begin{warning}
\textbf{An over-general predicate, recorded but not repaired.} \lean{RegularImpliesTrivialGraded W}
is stated for an arbitrary VOA $W$, whereas Theorem~3.1 assumes $W$ simple affine. For $W$ with
$V_1 = 0$ regularity is vacuous and the predicate would assert full graded rigidity, which is
morally false (e.g.\ for Virasoro $V(c,0)$, along $\partial_c$). This does not affect any result of this
chapter: the predicate is only ever \emph{proved}, never assumed, and only at affine $\lninesl$
quotients (\lean{L43_regularImpliesTrivialGraded}, \lean{simple_regularImpliesTrivialGraded}). A
reader should simply not cite \lean{isRigidGraded_of_split} as ``Theorem 3.1'' for general $W$.
\end{warning}

\subsection{The unconditional theorem at \texorpdfstring{$-4/3$}{-4/3}}

\begin{leancode}
theorem singular_central_polynomial (δ : ℂ) (v : Q) :
    let K : End ℂ (Q × Q) := kk • 1 + δ • tangentShift
    (24960 : ℂ) • K (v, 0) + (27360 : ℂ) • (K ^ 2) (v, 0) +
      (6480 : ℂ) • (K ^ 3) (v, 0) = (0, (-13440 * δ) • v) := by
\end{leancode}
\begin{leancode}
theorem L43_normalisedPoleVanishes : NormalisedPoleVanishes L43 := by
\end{leancode}
\begin{leancode}
theorem L43_graded_rigidity : Deform.IsRigidGraded L43 :=
  L43_isRigidGraded_of_normalisedPoleVanishes L43_normalisedPoleVanishes
\end{leancode}
(\leanpath{VOA/Examples/AffineSl2RigidityComplete.lean}; \lean{kk} is local notation for $-4/3$,
\lean{Q} for the carrier of \lean{L43}, \lean{tangentShift} is $(x_0,x_1)\mapsto(0,x_0)$ on $Q\times Q$, i.e.\
multiplication by $t$ on $Q\oplus tQ$.)

\begin{theorem}[Graded first-order rigidity of $L_{-4/3}(\lninesl)$; Linshaw--Qi, Theorem 5.1]\label{l09:thm:L43}
Every graded first-order deformation of $L_{-4/3}(\lninesl)$ is trivial.
\end{theorem}
\begin{physics}
$L_{-4/3}(\lninesl)$, central charge $-6$, is the chiral algebra associated with the $(A_1,D_3)$ Argyres--Douglas
theory (physics literature). Its null state at weight 3 encodes a relation among the currents that holds
only at $k=-4/3$; any attempt to move the level violates it at first order, with rate proportional to the
derivative $-13440$ of the null-norm polynomial $240k(3k+4)(9k+26)$. Hence no deformation.
\end{physics}


\begin{proof}[Proof on paper]
Normalise (Corollary~\ref{l09:cor:normalise}); it remains to show the normalised pole vanishes
(Piece~2), after which regular reconstruction (Section~\ref{l09:sec:uniform}) finishes. In pole form the
pole is $a^1_{(1)}b = \kappa(a,b)\delta\lnineone$ for a single scalar $\delta$, and the lifted currents
$\tilde J^a_n = J^a_n + tJ^{a,1}_n$ on $Q\oplus tQ$ satisfy the affine relations with central element
$K = k + \delta t$. Pair the singular word: the expression of $s$ in lifted currents, applied to the
vacuum, and then paired with the vacuum through the adjoint word, can be evaluated in two ways. By
normal ordering in the affine algebra with central element $K$ it equals
$240K(3K+4)(9K+26)$ applied to the vacuum --- the displayed cubic $24960K + 27360K^2 + 6480K^3$ ---
whose value at $K = -4/3 + \delta t$ is $0 + t(-13440\delta)$, because $k=-4/3$ is a simple root of the
cubic. On the other hand, since $s=0$ in $L_{-4/3}$, the lifted word lands in $tQ$, and the adjoint
relation makes its pairing vanish (\lean{actual_singular_pairing_obstruction}). Hence
$-13440\delta = 0$.
\end{proof}

The factor $(3K+4)$ in $240K(3K+4)(9K+26)$ is the reason the singular vector exists exactly at
$k=-4/3$; the \emph{simplicity} of that root is what makes the tangent coefficient nonzero. This is
the special case of the uniform mechanism of the next two sections, where ``simple root of an
explicit cubic'' is replaced by ``simple root of a Gram determinant''.

\section{The uniform machinery at every nonzero non-critical level}\label{l09:sec:uniform}

The files \leanpath{AffineSl2General*.lean} and \leanpath{AffineSl2UniformObstruction.lean} redo every
step of the $-4/3$ proof except the final singular-vector computation, \emph{for every level} $k$ with
$k\ne0$ and $k\ne-2$. They are stated under \lean{variable (k : ℂ) (hk : k + 2 ≠ 0) (hk0 : k ≠ 0)}. No
admissibility enters until Section~\ref{l09:sec:determinant}.

\subsection{The first-order extension as a module}

A first-order deformation $D$ of $W$ makes $W_R = W\oplus tW$ a vertex algebra over $R=\C[t]/(t^2)$. The
library models $W\oplus tW$ as the product $W\times W$, a pair $(x_0,x_1)$ standing for $x_0 + tx_1$. The
lifted modes and the full product are
\begin{equation*}
  \tilde u_{(n)}(x_0,x_1) = \bigl(u_{(n)}x_0,\ u^1_{(n)}x_0 + u_{(n)}x_1\bigr),\qquad
  (x_0,x_1)_{(n)}(y_0,y_1) = \bigl(x_{0(n)}y_0,\ x_{0(n)}y_1 + x_{1(n)}y_0 + x^1_{0(n)}y_0\bigr),
\end{equation*}
(\lean{liftedMode}, \lean{Deform.Reconstruction.product}) and multiplication by $t$ is
\lean{tangentShift}$\,(x_0,x_1) = (0,x_0)$.

\subsection{Step A: regular deformations are trivial (Theorem 3.1 at every level)}

\noindent File \leanpath{VOA/Examples/AffineSl2GeneralRegularReconstruction.lean}.
\begin{leancode}
theorem simple_regular_isTrivial (D : FirstOrderDeformation (simple k hk))
    (hD : D.IsGraded) (hreg : D.Regular) : D.IsTrivial := by
\end{leancode}
\noindent File \leanpath{VOA/Examples/AffineSl2GeneralRegularReconstruction.lean}.
\begin{leancode}
theorem simple_regularImpliesTrivialGraded : RegularImpliesTrivialGraded (simple k hk) :=
  (simple_regular_isTrivial k hk hk0)
\end{leancode}
The general reconstruction criterion behind it (\leanpath{VOA/Core/DeformationReconstruction.lean}):
\begin{leancode}
theorem isTrivial_of_closed_lift (S : Submodule ℂ (V × V))
    (hsimple : W.IsSimple) (hclosed : ModeClosed D S)
    (hsurj : Function.Surjective (reduction S))
    (hsep : (0, W.vacuum) ∉ S) (hvac : (W.vacuum, 0) ∈ S)
    (hgrade : GradeClosed W S) : D.IsTrivial := by
\end{leancode}

\begin{proof}[Proof on paper]
\emph{The criterion.} Suppose $S\subseteq W\oplus tW$ is closed under all deformed products, contains
$(\lnineone,0)$ but not $(0,\lnineone)$, surjects onto $W$ under $(x_0,x_1)\mapsto x_0$, and is closed
under homogeneous components. The kernel of the projection on $S$ is $\{(0,y)\}\cap S$; its second
coordinates form an ideal of $W$ (closure under products with lifts of all states), not containing
$\lnineone$, hence zero by simplicity. So the projection is an isomorphism $S\cong W$; writing its
inverse as $x\mapsto(x,\psi x)$, closure under products says exactly that $u^1_{(n)}v = (\delta(-\psi))$
\ldots, i.e.\ $D$ is the coboundary of a grade-preserving $\varphi$ with $\varphi\lnineone = 0$.
\emph{The carrier.} If $D$ is regular, the lifted currents satisfy the affine relations at the
\emph{same} level $k$ on $W\oplus tW$ (the deformed commutator has no correction to either the bracket
or the central term), and they annihilate $(\lnineone,0)$ in non-negative modes. By the universal
property of the vacuum module there is a map $\mathrm{ev}: V^k(\lninesl)\to W\oplus tW$ sending
$J^{a_1}_{-n_1}\cdots\lnineone$ to $\tilde J^{a_1}_{-n_1}\cdots(\lnineone,0)$. Take $S$ its image. Its first
coordinate is the quotient map $V^k\to L_k$, hence surjective; it respects energy, so its weight-zero
part is $\C(\lnineone,0)$, which excludes $(0,\lnineone)$; it is closed under current modes and, by the
iterate (associativity) identity, under all deformed products. This replaces the step in Linshaw--Qi
where they cite Lian's Theorem~4.11 to conclude $W_{(0)} = \C\lnineone$; the library's file records a
concern that the citation may be circular (Lian's category assumes one-dimensional weight-zero
space) and shows the step is obtained directly from the universal property of $V^k(\lnineg)$.
\end{proof}

\begin{warning}
\textbf{Provenance (recorded in the source).} Linshaw--Qi state that the proof of Theorem~3.1 in v2 was
contributed by an AI model, in a revision prompted by errors in v1, and it cites Lian for $W_{(0)} =
\C\lnineone$. The library does not rely on that proof: \lean{simple_regular_isTrivial} is proved
independently, with the universal-property argument above.
\end{warning}

\subsection{Step B: the single level parameter}

\noindent File \leanpath{VOA/Examples/AffineSl2GeneralDeformationPole.lean}.
\begin{leancode}
theorem simple_existsUnique_pole_parameter (D : FirstOrderDeformation (simple k hk))
    (hD : D.IsGraded) (hnorm : D.InPoleForm) :
    ∃! δ : ℂ, ∀ a b : (simple k hk).grading 1,
      D.poleVec a b = ((currentPairing k hk hk0) a b * δ) • (simple k hk).vacuum := by
\end{leancode}
\noindent File \leanpath{VOA/Examples/AffineSl2GeneralNormalizedRepresentation.lean}.
\begin{leancode}
theorem simple_exists_normalized_representation (D : FirstOrderDeformation (simple k hk))
    (hD : D.IsGraded) (hnorm : D.InPoleForm) :
    ∃ (δ : ℂ) (ρ : AffSl2 →ₗ⁅ℂ⁆ End ℂ ((WakModule k ⧸ I) × (WakModule k ⧸ I))),
      (∀ a b : (simple k hk).grading 1, D.poleVec a b = ((currentPairing k hk hk0) a b * δ) • (simple k hk).vacuum) ∧
      (∀ x n, ρ (ofLoop kappaForm kappaForm_isInvariant x n) =
        liftedMode D ((affineCurrent k hk hk0) x) n) ∧
      ρ (AffineLie.C kappaForm kappaForm_isInvariant) =
        k • (1 : End ℂ ((WakModule k ⧸ I) × (WakModule k ⧸ I))) + δ • tangentShift ∧
      ρ (ell0 kappaForm kappaForm_isInvariant) = liftedDegree (W := (simple k hk)) := by
\end{leancode}

\begin{dictionary}
\lean{currentPairing k hk hk0 a b} is $\kappa(a,b)$ computed in the proved basis $(e,h,f)$ of $V_1$ (entries
$\kappa(e,f)=\kappa(f,e) = 1$, $\kappa(h,h)=2$). \lean{∃! δ} is unique existence. So the first theorem says: for a
graded deformation in pole form, $a^1_{(1)}b = \delta\,\kappa(a,b)\lnineone$ for one scalar $\delta$. The second
packages the lifted currents as a Lie algebra homomorphism \lean{ρ} from $\lninehat$
(\lean{AffSl2}; \lean{→ₗ⁅ℂ⁆} is a Lie algebra morphism) to $\End(Q\times Q)$, with $Q$ the carrier of
$L_k$ and \lean{I} the maximal ideal; the central element acts as $k + \delta t$ and $\ell_0$ as the lifted
degree operator. \lean{affineCurrent} identifies $\lninesl$ (coordinate colours) with $V_1$.
\end{dictionary}

\begin{proof}[Proof on paper]
For graded $D$ in pole form, $a^1_{(1)}b\in V_0 = \C\lnineone$, giving a bilinear form $\beta$ on $V_1$
with $a^1_{(1)}b = \beta(a,b)\lnineone$. The $\ell=0$ case of the deformed Jacobi identity is the deformed
commutator formula
\begin{equation*}
  [\tilde a_m, \tilde b_n] = \widetilde{[a,b]}_{m+n} + m\bigl(k\kappa(a,b) + t\,\beta(a,b)\bigr)\delta_{m+n,0},
\end{equation*}
(\lean{poleForm_current_commutator}); its Jacobi identity in the Lie algebra makes $\beta$ invariant,
and every invariant form on $\lninesl$ is a multiple $\delta\kappa$ (\lean{exists_currentPole_sl2}).
Uniqueness: $\beta(e,f) = \delta$ and $\lnineone\neq0$.
\end{proof}

\noindent File \leanpath{VOA/Examples/AffineSl2GeneralDeformationPole.lean}.
\begin{leancode}
theorem simple_isRigidGraded_of_normalisedPoleVanishes
    (hvan : NormalisedPoleVanishes (simple k hk)) : IsRigidGraded (simple k hk) :=
  isRigidGraded_of_sl2Basis (simple k hk) (simple_sl2BasisOnWeightOne k hk hk0)
    (simple_regularImpliesTrivialGraded k hk hk0) hvan
\end{leancode}
So, at every nonzero non-critical level, graded rigidity of $L_k(\lninesl)$ is reduced to: \emph{the
level parameter $\delta$ of every normalised graded deformation vanishes.}

\subsection{Step C: word pairings are polynomials in the level}

For a word $w = ((a_1,n_1),\ldots,(a_r,n_r))$ write $J_w = J^{a_1}_{n_1}\cdots J^{a_r}_{n_r}$. In any
representation of $\lninehat$ in which the vacuum is killed by non-negative modes, the vacuum
coefficient of $J_w\lnineone$ is computed by moving annihilation operators to the right with
\eqref{l09:eq:affinebracket}; the result is a polynomial $P_w(C)$ in the central element. For the
lifted representation $C = k + \delta t$, so the vacuum pairing is $P_w(k+\delta t) = P_w(k) + t\,\delta P_w'(k)$.

\noindent File \leanpath{VOA/Examples/AffineSl2GeneralPairing.lean}.
\begin{leancode}
theorem simple_normalized_word_tangent (D : FirstOrderDeformation (simple k hk))
    (hD : D.IsGraded) (hnorm : D.InPoleForm) (δ : ℂ)
    (hδ : ∀ a b : (simple k hk).grading 1,
      D.poleVec a b = (currentPairing k hk hk0 a b * δ) • (simple k hk).vacuum)
    (w : List (Fin 3 × ℤ)) :
    (simpleForm k hk).form (simple k hk).vacuum
      (liftedCurrentWord k hk D w ((simple k hk).vacuum,0)).2 =
        (NormalOrdering.wordPolynomial w).derivative.eval k * δ := by
\end{leancode}

\begin{dictionary}
\lean{simpleForm k hk} is the canonical invariant bilinear form of $L_k$ with $(\lnineone|\lnineone)=1$
(exists because the currents are quasi-primary, \lean{simple_quasiPrimary}). \lean{liftedCurrentWord
k hk D w} is $\tilde J_w$ on $Q\times Q$ (leftmost letter acting last), applied to $(\lnineone,0)$; \lean{.2}
takes the $t$-coefficient. \lean{NormalOrdering.wordPolynomial w} is the universal normal-ordering
polynomial $P_w$, and \lean{.derivative.eval k} is $P_w'(k)$. The statement: the $t$-part of the
vacuum pairing of $\tilde J_w(\lnineone,0)$ is $\delta P'_w(k)$.
\end{dictionary}

\subsection{Step D: two relations give an obstruction}

\noindent File \leanpath{VOA/Examples/AffineSl2UniformObstruction.lean}.
\begin{leancode}
theorem graded_rigidity_of_relations {ι κ : Type*} [Fintype ι] [Fintype κ]
    (u : ι → List (Fin 3 × ℤ)) (c : ι → ℂ)
    (v : κ → List (Fin 3 × ℤ)) (d : κ → ℂ)
    (hc : ∑ i, c i • currentWord k hk (u i) (simple k hk).vacuum = 0)
    (hd : ∑ j, d j • currentWord k hk (v j) (simple k hk).vacuum = 0)
    (hp : (relationPairPolynomial u c v d).derivative.eval k ≠ 0) :
    IsRigidGraded (simple k hk) := by
\end{leancode}
\noindent File \leanpath{VOA/Examples/AffineSl2UniformObstruction.lean}.
\begin{leancode}
noncomputable def relationPairPolynomial {ι κ : Type*} [Fintype ι] [Fintype κ]
    (u : ι → List (Fin 3 × ℤ)) (c : ι → ℂ)
    (v : κ → List (Fin 3 × ℤ)) (d : κ → ℂ) : Polynomial ℂ :=
  ∑ i, ∑ j, (c i * d j) • wordGramPolynomial (u i) (v j)
\end{leancode}
\noindent File \leanpath{VOA/Examples/AffineSl2UniformObstruction.lean}.
\begin{leancode}
noncomputable def wordGramPolynomial (u v : List (Fin 3 × ℤ)) : Polynomial ℂ :=
  ((-1:ℂ)^u.length) • NormalOrdering.wordPolynomial (adjointIndices u ++ v)
\end{leancode}

\begin{theorem}[relation obstruction]\label{l09:thm:relations}
Let $k\ne0,-2$. Suppose $\sum_i c_i J_{u_i}\lnineone = 0$ and $\sum_j d_jJ_{v_j}\lnineone = 0$ in $L_k(\lninesl)$
(two \emph{relations}), and let $P(C) = \sum_{i,j}c_id_j\,G_{u_i,v_j}(C)$ be their Gram polynomial, where
$G_{u,v}(k) = (J_u\lnineone\,|\,J_v\lnineone)$. If $P'(k)\neq0$, then $L_k(\lninesl)$ is graded rigid.
\end{theorem}
\begin{physics}
A relation is a null-state equation among normal-ordered current products. Its Gram polynomial is the
two-point function of two such composite operators as a function of the level. If the level could drift,
$k\to k+\delta t$, the relation would have to persist, so this two-point function would have to stay zero;
its first-order change $\delta P'(k)$ must vanish. Transversality $P'(k)\neq0$ is the statement that the
null relation genuinely ``pins'' the level.
\end{physics}


\begin{dictionary}
\lean{currentWord k hk w} is $J_w$ acting on $L_k$, so \lean{hc}, \lean{hd} are the two relations;
\lean{Fintype ι} makes the sums finite. \lean{wordGramPolynomial u v} is the polynomial whose value
at $k$ is $(J_u\lnineone|J_v\lnineone)$ (proved as \lean{wordGramPolynomial_eval}): the adjoint of
$J^a_n$ for the invariant bilinear form is $-J^a_{-n}$ (\emph{same} colour, opposite mode,
\lean{adjointIndices}, one sign per letter). The two relations may differ (\lean{u} and \lean{v}
independent), which matters because the bilinear form pairs opposite $h_0$-weights. The polynomial
is constructed uniformly by normal ordering; \lean{hp} is its derivative at the level.
\end{dictionary}

\begin{proof}[Proof on paper]
Normalise $D$ and let $\delta$ be its level parameter. Lift the column relation:
$\tilde v := \sum_jd_j\tilde J_{v_j}(\lnineone,0)$ has first coordinate $\sum_jd_jJ_{v_j}\lnineone = 0$, so
$\tilde v = (0,y)$ is a pure tangent vector. Pair it against the row: the vacuum coefficient of
$\sum_ic_i\tilde J_{u_i}^{\dagger}\tilde v$ has $t$-part $\sum_ic_i(\lnineone\,|\,J_{u_i}^\dagger y) =
(\sum_ic_iJ_{u_i}\lnineone\,|\,y) = 0$ by the row relation (\lean{actual_word_relation_obstruction}). But by
Step~C the same quantity is the $t$-part of $P(k+\delta t)$, namely $\delta P'(k)$
(\lean{relationPair_tangent}). Hence $\delta P'(k) = 0$, so $\delta = 0$, and
\lean{simple_isRigidGraded_of_normalisedPoleVanishes} applies.
\end{proof}

Physically this is the precise version of the preview: the relations hold at level $k$; a
deformation that moves the level to $k+\delta t$ must preserve them; their Gram pairing changes at
rate $\delta P'(k)$; so a level change is forbidden whenever the relations are ``transverse'' to the
level ($P'(k)\neq0$). For the free boson there are no relations, and the level change is allowed (as
pure gauge); for $V^k(\lninesl)$ there are no relations, and the level change is a genuine deformation.

\section{The Kac determinant and a simple root}\label{l09:sec:determinant}

It remains to produce, at each negative admissible level, two relations with $P'(k)\ne0$. The library
does this without writing down any singular vector, by a determinant argument.

\subsection{From a simple determinant root to two relations}

\begin{lemma}[kernel pair]\label{l09:lem:kernelpair}
Let $M(X)$ be a square matrix of polynomials over an integral domain, and $z$ a simple root of
$\det M$: $\det M(z) = 0\neq(\det M)'(z)$. Then there are vectors $r,s$ with $r^{\mathsf T}M(z) = 0$,
$M(z)s = 0$ and $r^{\mathsf T}M'(z)s\neq0$.
\end{lemma}
\begin{physics}
Purely technical: it extracts, from the vanishing of the Kac--Kazhdan determinant (the norm-matrix of
descendant states) at a given level, two null states whose mutual two-point function changes at first order
when the level is moved. It is the linear-algebra step of the null-vector argument.
\end{physics}


\noindent File \leanpath{VOA/Core/PolynomialMatrixSimpleRoot.lean}.
\begin{leancode}
theorem kernel_pair_of_simple_root [NoZeroDivisors R]
    (M : Matrix ι ι R[X]) (z : R)
    (hz : M.det.eval z = 0) (hd : M.det.derivative.eval z ≠ 0) :
    ∃ r c : ι → R,
      (∀ j, ∑ i, r i * (M i j).eval z = 0) ∧
      (∀ i, ∑ j, (M i j).eval z * c j = 0) ∧
      (∑ i, ∑ j, (r i * c j) * (M i j).derivative.eval z) ≠ 0 := by
\end{leancode}

\begin{proof}[Proof on paper]
Let $A = \operatorname{adj}M$. From $AM = MA = \det M\cdot 1$ at $z$: $A(z)M(z) = M(z)A(z) = 0$. Differentiate
$MA = \det M\cdot1$ and multiply by $A(z)$ on both sides: $A(z)M'(z)A(z) = (\det M)'(z)A(z)$ (the
term with $A'$ is killed by $M(z)A(z)=0$). Since $(\det M)'(z)\neq0$, $A(z)\neq0$ (by Jacobi's formula
$(\det M)' = \operatorname{tr}(A M')$). Pick $(a,b)$ with $A(z)_{ab}\neq0$ and take $r$ = row $a$ of $A(z)$,
$s$ = column $b$; then $r^{\mathsf T}M'(z)s = (A(z)M'(z)A(z))_{ab} = (\det M)'(z)A(z)_{ab}\neq0$.
\end{proof}

Apply this to the Gram matrix of the invariant form on a weight space of $V^k(\lninesl)$. The left
kernel vector $r$ gives a combination of PBW states in the radical of the form --- hence in the maximal
ideal by Li's theorem --- hence a relation in $L_k$; similarly the right kernel vector. And
$r^{\mathsf T}M'(k)s$ is exactly the derivative of their Gram polynomial. One subtlety: the
\emph{Shapovalov} form of a Verma module is contravariant for the Chevalley anti-involution ($e_n^\dagger
= f_{-n}$), while the VOA's invariant bilinear form has $(J^a_n)^\dagger = -J^a_{-n}$. They differ by the
Chevalley automorphism $\theta$ ($e\mapsto -f$, $f\mapsto-e$, $h\mapsto-h$), which the library constructs on
$V^k$ (\lean{vacuumChevalley}) and shows preserves the maximal ideal (\lean{vacuumChevalley_ideal_iff});
one of the two kernel vectors is therefore twisted by $\theta$. The resulting statement:

\begin{leancode}
theorem graded_rigidity_of_vacuum_simple_root (hk0 : k≠0) (n : ℕ) (c : ℤ)
    (hz : (vacuumChevalleyPolynomial n c).det.eval k=0)
    (hd : (vacuumChevalleyPolynomial n c).det.derivative.eval k≠0) :
    Deform.IsRigidGraded (simple k hk) := by
\end{leancode}
(\leanpath{VOA/Examples/AffineSl2VacuumSimpleRoot.lean}, a Chapter~\ref{ch:L08} file.)

\begin{dictionary}
\lean{vacuumChevalleyPolynomial n c} is a square matrix of polynomials in one variable $X$ (the level),
indexed by \lean{JointIndex n (2*c)}: the PBW monomials of $V^X(\lninesl)$ of energy $n$ and
$h_0$-weight $2c$ (\lean{currentEnergy}, \lean{currentCharge}, with $e,f,h$ contributing $+2,-2,0$). Its
$(i,j)$ entry is the normal-ordering polynomial of the vacuum pairing of the $i$-th and $j$-th
monomials with the Chevalley adjoint convention. \lean{.det.eval k} and \lean{.det.derivative.eval k}
are $\det M(k)$ and $(\det M)'(k)$. The theorem: \emph{if the level is a simple root of one such
vacuum Gram determinant, $L_k(\lninesl)$ is graded rigid} (at any nonzero non-critical $k$).
\end{dictionary}

\subsection{Locating a simple root: the Kac determinant with the horizontal factor removed}

Which block, and why is the root simple? By \eqref{l09:eq:firstsing}, at $k+2 = p/q$ the first new
singular vector of $V^k$ has energy $n = q(p-1)$ and $h_0$-weight $2(p-1)$, so take $c = p-1$. The
Gram determinant of $V^k$ in that block is controlled by the Kac--Kazhdan determinant of the Verma
module $M(k,\lambda)$ of highest weight $(k,\lambda)$, as a polynomial in $\lambda$. At $\lambda = 0$ the Verma
determinant vanishes identically in $k$ because of the horizontal singular vector $f_0v$, which spans a
copy of $M(k,-2)$ inside $M(k,0)$; the vacuum module is the quotient by it. Extracting the first
nonvanishing coefficient in $\lambda$ gives the polynomial identity

\noindent File \leanpath{VOA/Examples/AffineSl2KacVacuumCoefficient.lean}.
\begin{leancode}
theorem affine_vacuum_polynomial_identity (n : ℕ) (c : ℤ)
    (hc : 0<2*(n : ℤ)-c) :
    (vacuumChevalleyPolynomial n c).det*affineMinusTwoPolynomial n (c+1)=
      Polynomial.C (vacuumDeterminantScale n c)*affineLeadingPolynomial n c*affineVacuumProduct n c := by
\end{leancode}
\begin{equation}\label{l09:eq:detidentity}
  \det M^{\mathrm{vac}}_{n,c}(k)\cdot \det M^{(-2)}_{n,c+1}(k)
  = \sigma_{n,c}\cdot \Lambda_{n,c}(k)\cdot \prod_{\beta\ne\alpha,\ r}\bigl(\text{real-root factor}_{\beta,r}(k)\bigr)^{\mathrm{mult}} ,
\end{equation}
where $M^{(-2)}$ is the Gram block of $M(k,-2)$, $\sigma\neq0$ is a combinatorial scale, $\Lambda(k)$ is a
product of oscillator factors that is nonzero at every non-critical level
(\lean{affineLeadingPolynomial_eval_ne_zero}), and the product (\lean{affineVacuumProduct}) runs over
the real-root labels of the Kac--Kazhdan formula \emph{other than the horizontal one}, each factor
being the affine-linear function of $k$ that vanishes exactly when \eqref{l09:eq:KK} holds, raised to
the dimension of the corresponding shifted weight space. This identity is a Chapter~\ref{ch:L08}
result (the Kac determinant is proved there from the screening/Wakimoto construction); this chapter
uses it as a proved theorem.

\noindent File \leanpath{VOA/Examples/AffineSl2AdmissibleKacLabels.lean}.
\begin{leancode}
theorem admissible_zero_labels (p q : ℕ) (hp : 2≤p) (hq : 0<q)
    (hcop : p.Coprime q) (i : (Fin 3×ℕ)×ℕ)
    (hi : i∈affineKacLabels (q*(p-1)) ((p : ℤ)-1))
    (hz : realRootWeight ((p : ℂ)/q) i=0) :
    i=((1,0),0) ∨ i=((0,q-1),p-2) := by
\end{leancode}
\noindent File \leanpath{VOA/Examples/AffineSl2AdmissibleVacuumRoot.lean}.
\begin{leancode}
theorem admissible_first_multiplicity (p q : ℕ) (hp : 2≤p) (hq : 0<q) :
    affineKacMultiplicity (q*(p-1)) ((p : ℤ)-1) ((0,q-1),p-2)=1 := by
\end{leancode}
\noindent File \leanpath{VOA/Examples/AffineSl2AdmissibleVacuumRoot.lean}.
\begin{leancode}
theorem affineMinusTwoPolynomial_admissible_ne_zero (p q : ℕ) (hp : 2≤p) (hq : 0<q)
    (hcop : p.Coprime q) :
    (affineMinusTwoPolynomial (q*(p-1)) (p : ℤ)).eval ((p : ℂ)/q-2)≠0 := by
\end{leancode}
\noindent File \leanpath{VOA/Examples/AffineSl2AdmissibleVacuumRoot.lean}.
\begin{leancode}
theorem vacuumChevalleyDeterminant_admissible_simple (p q : ℕ) (hp : 2≤p) (hq : 0<q)
    (hcop : p.Coprime q) :
    (vacuumChevalleyPolynomial (q*(p-1)) ((p : ℤ)-1)).det.eval ((p : ℂ)/q-2)=0 ∧
    (vacuumChevalleyPolynomial (q*(p-1)) ((p : ℤ)-1)).det.derivative.eval ((p : ℂ)/q-2)≠0 := by
\end{leancode}

\begin{dictionary}
A Kac label \lean{((a,m),r) : (Fin 3×ℕ)×ℕ} encodes a real root and an occurrence: colour \lean{a = 0}
is the family $(m+1)\delta-\alpha$, colour \lean{1} the family $m\delta+\alpha$ (with \lean{((1,0),0)} the
horizontal root $\alpha$ at its first occurrence), and \lean{r+1} the multiplicity $n$ in
\eqref{l09:eq:KK}. \lean{realRootWeight t i} is the affine-linear function of the shifted level
$t = k+2$ that vanishes on the Kac--Kazhdan hyperplane. \lean{admissible_zero_labels} says: in the
block of energy $q(p-1)$ and charge $p-1$, at $t = p/q$, \emph{only two} labels vanish --- the
horizontal one (removed in \lean{affineVacuumProduct}) and \lean{((0,q-1),p-2)}, i.e.\ $\beta = q\delta-\alpha$
with $n = p-1$, which is \eqref{l09:eq:firstsing}. \lean{admissible_first_multiplicity} says that
factor occurs to the first power (the shifted weight space is the vacuum line only).
\lean{affineMinusTwoPolynomial_admissible_ne_zero} says $\det M^{(-2)}\neq0$ at the admissible level.
\lean{(p : ℂ)/q-2} is $k$. Coprimality is used in \lean{admissible_zero_labels} (through the
resonance arithmetic) and in the nonvanishing of the $M(k,-2)$ determinant.
\end{dictionary}

\begin{proof}[Proof on paper that the root is simple]
Evaluate \eqref{l09:eq:detidentity} at $k = p/q-2$. The right side has exactly one vanishing factor,
the linear factor $-q\bigl(X - (p/q-2)\bigr)$ to the first power (\lean{admissible_first_root_polynomial});
everything else on the right is nonzero, and $\det M^{(-2)}(k)\neq0$. Hence $\det M^{\mathrm{vac}}(k) = 0$.
Differentiate the identity and evaluate at $k$: the left side gives $(\det M^{\mathrm{vac}})'(k)\det M^{(-2)}(k)$
(the other term has the factor $\det M^{\mathrm{vac}}(k)=0$); the right side gives
$\sigma\Lambda(k)\cdot(-q)\cdot(\text{other factors at }k)\neq0$ (the terms differentiating the other factors are
multiplied by the vanishing linear factor). So $(\det M^{\mathrm{vac}})'(k)\neq0$.
\end{proof}

\begin{remark}
Nothing here asserts that the kernel is spanned by ``the'' singular vector, nor that the maximal ideal
is generated by it; and no character formula is used. The determinant only certifies that \emph{some}
radical vector exists in the block, transversally to the level. This is exactly what the deformation
argument needs, and it is why the uniform proof is shorter than a proof via Adamovi\'c--Milas.
\end{remark}

\section{The headline theorem}\label{l09:sec:headline}

\begin{leancode}
theorem admissible_level_noncritical (p q : ℕ) (hp : 2≤p) (hq : 0<q) :
    (p : ℂ)/q-2+2≠0 := by
\end{leancode}
\begin{leancode}
theorem negative_admissible_level_ne_zero (p q : ℕ) (hq : 0<q) (hneg : p<2*q) :
    (p : ℂ)/q-2≠0 :=
\end{leancode}
\begin{leancode}
theorem negative_admissible_graded_rigidity (p q : ℕ) (hp : 2≤p) (hq : 0<q)
    (hcop : p.Coprime q) (hneg : p<2*q) :
    Deform.IsRigidGraded (simple ((p : ℂ)/q-2) (admissible_level_noncritical p q hp hq)) :=
\end{leancode}
(\leanpath{VOA/Examples/AffineSl2AdmissibleRigidity.lean}, namespace \lean{VOA.AffineSl2GeneralRigidity}.
The proof is shown in full because it is informative: it is the composition of the simple-root theorem
of Section~\ref{l09:sec:determinant} with the rigidity criterion.)

\begin{theorem}[graded first-order rigidity at negative admissible levels]\label{l09:thm:headline}
Let $p,q\in\N$ with $p\ge2$, $q\ge1$, $\gcd(p,q)=1$ and $p<2q$, and let $k = p/q-2$. Then every graded
first-order deformation of the simple affine vertex operator algebra $L_k(\lninesl)$ is trivial: for
every $Y_1$ satisfying the order-$t$ vertex-algebra axioms and the weight rule, there is a linear
$\varphi:L_k\to L_k$ preserving each weight space, with $\varphi\lnineone = 0$ and $Y_1 = \delta\varphi$.
\end{theorem}
\begin{physics}
For every negative fractional admissible level, the chiral algebra of the fractional-level $SU(2)$ WZW
model has no chiral moduli: its OPE data are isolated. Equivalently, one cannot continuously connect two such
theories at different admissible levels through a family of chiral algebras, which is consistent with the
levels being discrete (rational) labels fixed by modularity. The universal $V^k$, by contrast, has the level
as a modulus. Rigidity here is not explained by unitarity or rationality ($L_k$ is neither), but by the null
vector of energy $q(p-1)$.
\end{physics}


\begin{dictionary}
\textbf{The variables.} \lean{p q : ℕ} are natural numbers. \lean{hp : 2≤p} is $p\ge2$ (the
admissibility bound $p\ge h^\vee$); \lean{hq : 0<q} is $q\ge1$; \lean{hcop : p.Coprime q} is
$\gcd(p,q)=1$; \lean{hneg : p<2*q} is $p<2q$, equivalently $k = p/q-2 < 0$. These are exactly
``$k$ negative admissible'' (Definition~\ref{l09:def:admissible}); no further hypothesis is present.
\textbf{The level.} \lean{(p : ℂ)/q-2} is $k = p/q-2$ computed in $\C$ (the cast \lean{(p : ℂ)} and the
coercion of \lean{q} to $\C$). \lean{admissible_level_noncritical p q hp hq} is the proof of $k+2\ne0$
required to form the algebra; it is a proof term, and by proof irrelevance its choice cannot
matter. \textbf{The object.} \lean{simple k hk} is
\lean{QuotVOA.maxQuotientVOA (affineSl2VOA_closed k hk)}: the quotient of the \emph{universal} affine
VOA $V^k(\lninesl)$ (Sugawara conformal vector, $c = 3k/(k+2)$) by its maximal proper graded ideal,
i.e.\ the \emph{simple} quotient $L_k(\lninesl)$ (Definition~\ref{l09:def:Lk}). It is proved simple and
nonzero (\lean{simple_isSimple}, \lean{simple_vacuum_ne_zero}) and its weight-one space is $\lninesl$
(\lean{simple_sl2BasisOnWeightOne}). \textbf{The conclusion.} \lean{Deform.IsRigidGraded} is
Definition~\ref{l09:def:rigid}: for all \lean{D : FirstOrderDeformation (simple k hk)} with
\lean{D.IsGraded}, \lean{D.IsTrivial}. \textbf{Quantifier order.} $\forall p\,q$, (hypotheses)
$\Rightarrow \forall D$ graded, $\exists\varphi$: the coboundary $\varphi$ may depend on $D$, as it must.
\textbf{Same, stronger or weaker than the literature?} At $(p,q) = (2,3)$ it is Linshaw--Qi
Theorem~5.1, in a form at least as strong as theirs (Section~\ref{l09:sec:deform}: the library allows
deformations of the translation operator, theirs does not; the trivial ones agree). For all other
negative admissible levels it is new; it confirms, for $\lninesl$ and $k<0$, Linshaw--Qi's conjecture
that $L_k(\lnineg)\neq V^k(\lnineg)$ implies rigidity. It is \emph{about the simple quotient}, not the
universal algebra, and that matters: $V^k(\lninesl)$ is not rigid (the level is a genuine deformation
parameter), so a theorem about $V^k$ would have been false.
\textbf{How could it have been vacuous, and why it is not.} (a) \emph{Zero algebra:} excluded,
$\lnineone\neq0$, proved without \lean{VOA.pbw}. (b) \emph{Trivially rigid algebra:} $L_k$ is the genuine
simple quotient with $V_1\cong\lninesl$ and the non-zero level $k$, so the obvious candidate deformation
(change of level) is a real constraint; the analogous statement fails for the free boson
(\lean{not_normalisedPoleVanishes_heisenberg}) and for $V^k$ (mathematically). (c) \emph{Unsatisfiable
hypotheses:} the hypotheses are numerical and satisfied, e.g.\ by $(2,3)$, $(3,2)$, $(2,5)$; the
library also states $k=-4/3$ directly (\lean{L43_graded_rigidity}). (d) \emph{An empty deformation
class:} \lean{FirstOrderDeformation} is the full space of cocycles, with finite Jacobi sums. (e)
\emph{A hidden premise:} \texttt{\#print axioms} shows only Lean's standard axioms, and the statement
has no structure-valued hypothesis (no \lean{BottomProjection}-type datum). The general
non-vacuity caveat recorded in Section~\ref{l09:sec:deform} is now closed: the library exhibits $V(c,0)$ with a
non-trivial graded deformation (\lean{VOA.}\allowbreak\lean{VirPoly.}\allowbreak\lean{not_}\allowbreak\lean{isRigidGraded_}\allowbreak\lean{virasoroVOA}).
\end{dictionary}

\begin{status}
\lean{VOA.AffineSl2GeneralRigidity.negative_admissible_graded_rigidity}: \textbf{graded first-order
rigidity, unconditional} (hypotheses are only the numerical admissibility conditions; non-vacuity of the
notion machine-checked via \lean{VOA.}\allowbreak\lean{VirPoly.}\allowbreak\lean{not_}\allowbreak\lean{isRigidGraded_}\allowbreak\lean{virasoroVOA}). \texttt{\#print axioms}:
\texttt{[propext, Classical.choice, Quot.sound]}; no \lean{VOA.pbw}. Re-run for this chapter.
\end{status}

\begin{proof}[Proof on paper, assembled]
Let $D$ be graded. (1) $V_1\cong\lninesl$ since $k\neq0$; normalise $D$ by Whitehead
(Section~\ref{l09:sec:whitehead}). (2) The normalised $D$ has a unique level parameter $\delta$, and the
lifted currents represent $\lninehat$ at level $k+\delta t$ (Step B). (3) At $k+2 = p/q$, the vacuum Gram
determinant in energy $q(p-1)$, $h_0$-weight $2(p-1)$, has a simple zero at $k$
(Section~\ref{l09:sec:determinant}); its adjugate gives two relations in $L_k$ whose Gram polynomial has
nonzero derivative. (4) The relations force $\delta P'(k) = 0$, so $\delta=0$ (Step D). (5) Then the
normalised deformation is regular, hence trivial by reconstruction (Step A), and $D$ is trivial.
\end{proof}

\subsection{Every admissible level: the hypothesis \texorpdfstring{$p<2q$}{p<2q} is used only for \texorpdfstring{$k\ne0$}{k != 0}}\label{l09:sec:alladmissible}

\begin{sloppypar}In the proof of the headline the hypothesis \lean{hneg : p<2*q} enters only to show $k\neq0$; the
simple-root theorem \lean{vacuumChevalleyDeterminant_admissible_simple} needs only $p\ge2$, $q\ge1$, $\gcd(p,q)=1$.
Since $k = 0$ means $p = 2q$, which with coprimality forces $(p,q) = (2,1)$, the same argument proves rigidity at
\emph{every admissible level except $k=0$}, including the positive integral levels (Linshaw--Qi's Theorem~4.1 for
$\lninesl$) and the positive fractional ones (e.g.\ $k = 1/2$). The library states this form, and the headline is
its specialization to $p<2q$.\end{sloppypar}

\begin{theorem}[Graded first-order rigidity at every admissible level]\label{l09:thm:alladmissible}
Let $p,q$ be coprime positive integers with $p\ge2$ and $p\neq2q$, and put $k=p/q-2$. Then every graded first-order
deformation of the simple affine VOA $L_k(\lninesl)$ is trivial.
\end{theorem}

\noindent File \leanpath{VOA/\allowbreak Examples/\allowbreak AffineSl2AdmissibleRigidity.lean}.
\begin{leancode}
theorem admissible_graded_rigidity (p q : ℕ) (hp : 2≤p) (hq : 0<q)
    (hcop : p.Coprime q) (hne : p ≠ 2*q) :
    Deform.IsRigidGraded (simple ((p : ℂ)/q-2) (admissible_level_noncritical p q hp hq)) := by
\end{leancode}

\begin{dictionary}
The statement is that of Theorem~\ref{l09:thm:headline} with \lean{hneg : p<2*q} replaced by \lean{hne : p ≠ 2*q},
which excludes only $k=0$, where $L_0(\lninesl)=\C$. The headline
\lean{negative_admissible_graded_rigidity} is proved from it, discharging \lean{hne} from \lean{hneg} by \lean{omega}.
\end{dictionary}

\begin{status}
Proved; \texttt{\#print axioms}: \lean{[propext, Classical.choice, Quot.sound]} (re-run on the release build of
3~November~2026). No hypothesis beyond the four numerical conditions.
\end{status}

\begin{physics}
At positive integral $k$ this is the chiral algebra of the unitary $SU(2)$ WZW model; at fractional admissible $k$ it is
a non-unitary CFT. Either way the current OPE admits no infinitesimal deformation other than field redefinitions.
\end{physics}

\begin{remark}\label{l09:rem:extension}
Combined with the Gorelik--Kac criterion (for $\lninesl$, $V^k$ is non-simple at non-critical $k$ exactly when
$k+2\in\Q_{>0}$ has numerator at least $2$, i.e.\ $k$ admissible --- a literature fact, not formalised here), this
means that for $\lninesl$ Linshaw--Qi's conjecture ``$L_k\neq V^k\Rightarrow$ rigid'' holds at every non-critical
level except $k=0$, where $L_0 = \C$ and the conjecture is trivial. This combines a library theorem with a literature
fact; the combination is not a library theorem.
\end{remark}

\section{The boundary-admissible \texorpdfstring{$C_2$}{C2} programme (conditional)}\label{l09:sec:boundary}

The files \leanpath{AffineSl2Boundary*.lean} assigned to this chapter are not about deformations. They
concern the \emph{Zhu $C_2$ algebra} (Poisson algebra) $R_V = V/C_2(V)$, $C_2(V) =
\operatorname{span}\{u_{(-2)}v\}$, of the boundary-admissible algebras $L_k(\lninesl)$,
$k = -2 + 2/(2n+1)$, which governs their associated variety and quasi-lisse property (on paper; associated varieties are not formalised,
see Chapters~\ref{ch:L11} and~\ref{ch:L16}).

\begin{warning}[title={Superseded: now proved unconditionally for every boundary level}]
This section records the conditional programme as it stood when the chapter was written. The Arakawa--Dai campaign
has since proved, for every $m\ge1$ and $k=-2+2/(2m+1)$ and with no hypothesis, the exact presentation of the $C_2$
algebra of the actual simple VOA (\lean{AffineSl2BoundaryClassicalFreeness.baseC2AlgEquiv}), classical freeness and the
full vacuum character, bundled as the single theorem
\lean{VOA.}\allowbreak\lean{AffineSl2BoundaryClassicalFreeness.}\allowbreak\lean{arakawaDai_theorem_1_1}
(\leanpath{VOA/\allowbreak Examples/\allowbreak AffineSl2BoundaryTheorem11.lean}); see Chapter~\ref{ch:L20}. The conditional statements below remain true and proved, but they are no
longer the state of the art.
\end{warning}

For the universal algebra $R_{V^k(\lninesl)} = \operatorname{Sym}(\lninesl)
= \C[e,h,f]$; with $\mathrm{Cas} = h^2+4ef$ the expected answer (Arakawa--Jiang--Moreau, and the
classical literature) is
\begin{equation*}
  R_{L_k(\lninesl)} \cong \C[e,h,f]\big/\bigl(\mathrm{Cas}^n e,\ \mathrm{Cas}^n h,\ \mathrm{Cas}^n f\bigr),\qquad
  k = -2+\tfrac{2}{2n+1}.
\end{equation*}
These results are \textbf{conditional}: the library proves the presentation \emph{from} a named
representation-theoretic input (that the maximal ideal is generated by the singular vector, the
Adamovi\'c--Milas theorem), which it deliberately does not install as an axiom.

\noindent File \leanpath{VOA/Examples/AffineSl2BoundaryC2.lean}.
\begin{leancode}
def boundaryIdeal (n : ℕ) : Ideal SymSl2 :=
  Ideal.span
    ({casSym ^ n * eSym, casSym ^ n * hSym, casSym ^ n * fSym} : Set SymSl2)
\end{leancode}
\noindent File \leanpath{VOA/Commutative/Presentation.lean}.
\begin{leancode}
def IsExactPresentation (I : Ideal A) (φ : A →ₐ[R] B) : Prop :=
  RingHom.ker φ = I
\end{leancode}
\noindent File \leanpath{VOA/Examples/AffineSl2BoundaryC2Wakimoto.lean}.
\begin{leancode}
theorem exactPresentation43_of_maximalGeneration
    (hmax : QuotVOA.maxGradedIdeal
        (AffineVOA.affineSl2VOA_closed (-4 / 3 : ℂ) AffineVOA.Inst.admLevel_noncritical) =
      SingVec.idealGen
        (AffineVOA.affineSl2VOA_closed (-4 / 3 : ℂ) AffineVOA.Inst.admLevel_noncritical)
        Sl2Sing.sl2Singular) :
    Presentation.IsExactPresentation (AffineSl2BoundaryC2.boundaryIdeal 1)
      (AffineSl2BoundaryC2.quotientSymbol (-4 / 3 : ℂ) AffineVOA.Inst.admLevel_noncritical
        (QuotVOA.maxGradedIdeal
          (AffineVOA.affineSl2VOA_closed (-4 / 3 : ℂ) AffineVOA.Inst.admLevel_noncritical))
        (QuotVOA.maxGradedIdeal_isIdeal _)
        (QuotVOA.maxGradedIdeal_isGradedIdeal _)) :=
  exactPresentation_of_classicalData (classicalData43 hmax)
\end{leancode}
\noindent File \leanpath{VOA/Examples/AffineSl2BoundaryMaximalIdeal.lean}.
\begin{leancode}
def OnlyVacuumAffinePrimitives (s : WakModule k) : Prop :=
  ∀ v : WakModule k ⧸ idealGen (W k hk) s,
    SingularQuotientAffinePrimitive k hk s v →
      v ∈ ℂ ∙ (quotientBySingular (W k hk) s).vacuum
\end{leancode}

\begin{dictionary}
\lean{SymSl2} is \lean{MvPolynomial (Fin 3) ℂ}, the polynomial ring $\C[e,f,h]$, and \lean{casSym} is
$h^2+4ef$. \lean{boundaryIdeal n} is $(\mathrm{Cas}^ne,\mathrm{Cas}^nh,\mathrm{Cas}^nf)$ (for $n=1$ it equals the
earlier \lean{casIdeal}, \lean{boundaryIdeal_one}). \lean{quotientSymbol k hk I …} is the composite
$\C[e,f,h] = R_{V^k}\to R_{V^k/I}$, which is surjective (\lean{quotientSymbol_surjective}).
\lean{IsExactPresentation I φ} is $\ker\varphi = I$; together with surjectivity it gives
$R_{L_k}\cong\C[e,f,h]/I$ (\lean{boundarySimpleQuotientEquiv_of_singularC2Image}).
\lean{SingVec.idealGen W s} is the VOA ideal generated by $s$. So
\lean{exactPresentation43_of_maximalGeneration} reads: \emph{if} the maximal ideal of $V^{-4/3}$ is
generated by the weight-3 singular vector, \emph{then} $R_{L_{-4/3}}\cong\C[e,f,h]/(\mathrm{Cas}\,e,\mathrm{Cas}\,h,\mathrm{Cas}\,f)$.
The same holds at $k=-8/5$ (\lean{exactPresentation85_of_maximalGeneration}, $n=2$, with the
weight-5 singular vector). \lean{OnlyVacuumAffinePrimitives s}: in $V^k/(s)$, every affine primitive
vector (nonzero, killed by all positive modes and by $e_0$) lies on the vacuum line. The library proves
this is \emph{equivalent} to maximal generation (\lean{affineVacuumMaximalSubmoduleGenerated_of_onlyVacuumAffinePrimitives},
\lean{exactPresentation43_of_onlyVacuumAffinePrimitives}), so the sharpest remaining external input is
that primitive-uniqueness statement; the existence of primitives in every nonzero proper current
submodule is proved internally. The general-$n$ interface (\lean{AffineSl2BoundarySingularFamily}:
structures \lean{Candidate}, \lean{PrimitiveCriterionData}, \lean{AdamovicMilasData},
\lean{ClassicalData}) carries the singular-vector data as structure fields; concrete instances exist only
for $n = 1,2$ ($p = 3,5$).
\end{dictionary}

\begin{proof}[Proof on paper of the conditional presentation]
Three steps, all proved. (1) For a graded ideal $I$, $\ker(R_W\to R_{W/I})$ is the image of $I$ in $R_W$
(\lean{ker_A₂quot_eq_idealImage}). (2) Lower bound: the image of an ideal is a Poisson ideal; the
leading symbol of $s$ is a nonzero multiple of $e\,\mathrm{Cas}^n$ (computed for $n=1,2$), and applying the
horizontal lowering bracket twice produces $\mathrm{Cas}^nh$ and $\mathrm{Cas}^nf$, so the boundary ideal is
contained in the kernel. (3) Upper bound: for a vector killed by all positive current modes, the
ideal it generates is spanned by \emph{non-positive} current words applied to it
(\lean{singularCurrentPBWSpanning_of_SingP}); modulo $C_2$, modes $\le-2$ vanish, $-1$ is multiplication
by a symbol and $0$ is the horizontal Poisson bracket, which preserves the boundary ideal. Injectivity
of $\C[e,f,h]\to R_{V^k}$ (the universal PBW theorem, proved via Wakimoto) converts equality of images
into the exact kernel.
\end{proof}

\begin{status}
\lean{exactPresentation43_of_maximalGeneration}, \lean{exactPresentation85_of_maximalGeneration},
and the \lean{_of_onlyVacuumAffinePrimitives} and \lean{_of_currentIrreducible} variants:
\textbf{conditional} on the displayed hypothesis (maximal generation / primitive uniqueness /
irreducibility of the singular quotient), which is a published theorem (Adamovi\'c--Milas, Theorem
3.1.4) not formalised here. \texttt{\#print axioms} for the $-4/3$ statement (re-run):
\texttt{[propext, Classical.choice, Quot.sound]}. Nothing in this section is used by the rigidity
theorems.
\end{status}

\begin{warning}
\textbf{What the free-field route does not give.} An Adamovi\'c inverse-DS/free-field realisation shows
that the candidate singular vector lies in the maximal ideal
(\lean{mem_maxGradedIdeal_of_adamovic_freeField}), but \emph{not} that it generates it: the descended
map $V^k/(s)\to$ (free fields) is injective exactly when no further subsingular vectors exist, which a
finite check of $s\mapsto0$ cannot see (\lean{onlyVacuumAffinePrimitives_iff_strictDescent},
\lean{strictDescent_iff_maximalGeneration}). Treating ``$s$ vanishes in the free-field realisation'' as
``$s$ generates the maximal ideal'' would be an over-claim.
\end{warning}

\section{A cited hypothesis: admissible partitions}\label{l09:sec:partitions}

\leanpath{VOA/Examples/AdmissiblePartitions.lean} is assigned here by name but belongs to the
graded-unitarity programme for the Virasoro boundary minimal models $L(c_{2,2n+1},0)$
(Chapter~\ref{ch:L19}). Its purpose is to make one imported combinatorial input visible as a named
proposition, rather than proving anything:

\begin{leancode}
def AdmissibleBasisHyp (c : ℂ) (a : ℕ → ℕ → ℕ) : Prop :=
  ∀ N p : ℕ, finrank ℂ (goodFil c p (N : ℤ)) = ∑ r ∈ Finset.range (p + 1), a N r
\end{leancode}

\begin{dictionary}
\lean{goodFil c p N} is the $p$-th step of the length filtration on the weight-$N$ space of the simple
Virasoro quotient, and \lean{finrank} its dimension; \lean{a N r} is a \emph{parameter} intended to be
the number of Salazar's admissible partitions of $N$ of length $r$ (arXiv:2405.07126v3, Theorem~0.1).
The predicate is the Andrews--Gordon/Salazar refined character. It is \textbf{not proved} and is carried
as a hypothesis by every statement that uses it; \lean{a} is deliberately not yet a definition. The
file also records (as a docstring, numerically checked at $c=-22/5$, $N=6$) that the companion
\lean{FlagBasisHyp} must order the basis compatibly with length: reversing the order flips the Gram
pivot signs. None of this is related to affine rigidity.
\end{dictionary}

\section{Verification record}\label{l09:sec:verification}

All axiom checks were run for this chapter on a scratch file importing the relevant modules of the
already-built tree, with the command (never a bare \texttt{lake env lean}, which can drop
\texttt{maxSynthPendingDepth}):

\begin{plaincode}
export PATH="$HOME/.elan/bin:$PATH"
lake env lean -DmaxSynthPendingDepth=3 <scratchpad>/L09/Ax.lean
\end{plaincode}

with output (one line per declaration, abbreviated by namespace):

\begin{plaincode}
'VOA.AffineSl2GeneralRigidity.negative_admissible_graded_rigidity' depends on axioms: [propext, Classical.choice, Quot.sound]
'VOA.LQ43.L43_graded_rigidity' depends on axioms: [propext, Classical.choice, Quot.sound]
'VOA.LQ43.L43_isRigidGraded' depends on axioms: [propext, Classical.choice, Quot.sound]
'VOA.LQ43.isRigidGraded_of_split' depends on axioms: [propext, Classical.choice, Quot.sound]
'VOA.LQ43.CorrectedSystem.c_eq_zero' depends on axioms: [propext, Classical.choice, Quot.sound]
'VOA.LQ43.lqVector_eq_sl2Singular' depends on axioms: [propext, Classical.choice, Quot.sound]
'VOA.LQ43.maxGradedIdeal_ne_bot' depends on axioms: [propext, Classical.choice, Quot.sound]
'VOA.LQ43.L43_isSimple' depends on axioms: [propext, Classical.choice, Quot.sound]
'VOA.VertexOperatorAlgebra.weightOneH2Vanishes_of_sl2Basis' depends on axioms: [propext, Classical.choice, Quot.sound]
'VOA.Deform.weightOneCochain_mem_twoCocycle' depends on axioms: [propext, Classical.choice, Quot.sound]
'VOA.Sl2Whitehead.whitehead2_of_sl2Basis' depends on axioms: [propext, Classical.choice, Quot.sound]
'VOA.AffineSl2GeneralRigidity.graded_rigidity_of_vacuum_simple_root' depends on axioms: [propext, Classical.choice, Quot.sound]
'VOA.AffineSl2GeneralRigidity.simple_vacuum_ne_zero' depends on axioms: [propext, Classical.choice, Quot.sound]
'VOA.AffineSl2GeneralRigidity.simple_sl2BasisOnWeightOne' depends on axioms: [propext, Classical.choice, Quot.sound]
'VOA.AffineSl2GeneralRigidity.simple_regularImpliesTrivialGraded' depends on axioms: [propext, Classical.choice, Quot.sound]
'VOA.Wak.Charged.vacuumChevalleyDeterminant_admissible_simple' depends on axioms: [propext, Classical.choice, Quot.sound]
\end{plaincode}

and from a second scratch file (importing \leanpath{AffineSl2BoundaryC2Wakimoto},
\leanpath{AffineSl2RigidityCollapse}, \leanpath{AffineSl2AdmissibleRigidity}):

\begin{plaincode}
'VOA.AffineSl2BoundaryC2Wakimoto.exactPresentation43_of_maximalGeneration' depends on axioms: [propext, Classical.choice, Quot.sound]
'VOA.AffineSl2BoundaryC2Wakimoto.exactPresentation85_of_maximalGeneration' depends on axioms: [propext, Classical.choice, Quot.sound]
'VOA.AffineSl2RigidityCollapse.not_gradedImpliesRegular_L43' depends on axioms: [propext, Classical.choice, Quot.sound]
'VOA.AffineSl2RigidityCollapse.not_gradedImpliesPoleZero_L43' depends on axioms: [propext, Classical.choice, Quot.sound]
'VOA.AffineSl2RigidityCollapse.not_normalisedPoleVanishes_heisenberg' depends on axioms: [propext, Classical.choice, Quot.sound]
'VOA.AffineSl2RigidityCollapse.no_uniform_derivation_piece2' depends on axioms: [propext, Classical.choice, Quot.sound]
'VOA.AffineSl2RigidityCollapse.not_trivial_satisfies_piece2' depends on axioms: [propext, Classical.choice, Quot.sound]
'VOA.PolynomialMatrix.kernel_pair_of_simple_root' depends on axioms: [propext, Classical.choice, Quot.sound]
'VOA.AffineSl2GeneralRigidity.simple_regular_isTrivial' depends on axioms: [propext, Classical.choice, Quot.sound]
'VOA.AffineSl2GeneralRigidity.simple_isSimple' depends on axioms: [propext, Classical.choice, Quot.sound]
'VOA.AffineSl2GeneralRigidity.graded_rigidity_of_relations' depends on axioms: [propext, Classical.choice, Quot.sound]
'VOA.AffineSl2GeneralRigidity.simple_existsUnique_pole_parameter' depends on axioms: [propext, Classical.choice, Quot.sound]
'VOA.Wak.Charged.admissible_zero_labels' depends on axioms: [propext, Classical.choice, Quot.sound]
'VOA.Wak.Charged.affine_vacuum_polynomial_identity' depends on axioms: [propext, Classical.choice, Quot.sound]
\end{plaincode}

(Lean prints long lists over several lines; they are joined here.) The all-admissible theorem (Theorem~\ref{l09:thm:alladmissible}) gives the same three axioms. No full
\texttt{lake build} was run for this chapter (the tree was already built at the documented commit);
the project's build record for the commit that completed the mathematics reports an
aggregate build of 9{,}989 jobs and a full axiom audit of 3{,}310 reports, and the commit message of
the documented commit reports an aggregate build of 10{,}343 jobs with \texttt{check\_repository.py} and
\texttt{run\_axiom\_audits.py} passing. Those are the source's claims; the axiom lists above are what
this chapter's author re-confirmed.

\textbf{Re-run 2026-10-27.} At a later commit (branch \texttt{deligne-lean}; not re-run for this edition), with the same command on a
scratch file, \texttt{\#print axioms} gave exactly \texttt{[propext, Classical.choice, Quot.sound]} for
\lean{VOA.LQ43.L43_graded_rigidity}, \lean{VOA.AffineSl2GeneralRigidity.negative_admissible_graded_rigidity},
\lean{VOA.VirCDef.cmode_conformal_of_isCoboundary}, \lean{VOA.VirCDef.not_isTrivial_of_omega_witness},
\lean{VOA.VirCDef.not_isTrivial_of_isCDeformation}, \lean{VOA.VirCDef.not_isRigidGraded_of_isCDeformation} and
\lean{VOA.NonzeroDeform.not_isRigid_virasoroVOA}. No aggregate \texttt{lake build} was run.

Release audit 2026-10-27: \texttt{Notes/Textbook/referee/RELEASE\_AUDIT\_RIGIDITY.md}. Its verdict on
\lean{VOA.LQ43.L43_graded_rigidity} is ``announceable with caveats'' (graded first-order only; non-vacuity
argued on paper only). The non-vacuity caveat was removed afterwards by \lean{VOA.}\allowbreak\lean{VirPoly.}\allowbreak\lean{not_}\allowbreak\lean{isRigidGraded_}\allowbreak\lean{virasoroVOA}
(release follow-ups audit \texttt{Notes/Textbook/referee/RELEASE\_AUDIT\_FOLLOWUPS.md}, 2026-10-28, which reports
standard axioms for it; not re-run for this chapter). The audit also confirms that the comparison with Linshaw--Qi in Section~\ref{l09:sec:deform}
is correct: the Lean notion implies theirs.

\section{Referee's checklist}\label{l09:sec:checklist}

``Uncond.''\ means the hypotheses listed are ordinary mathematical hypotheses of the statement (no
unestablished datum); $\ast$ marks a \texttt{\#print axioms} run for this chapter with result
\texttt{[propext, Classical.choice, Quot.sound]}. No entry uses \lean{VOA.pbw}. Namespaces:
\lean{GR} = \lean{VOA.AffineSl2GeneralRigidity}, \lean{LQ} = \lean{VOA.LQ43},
\lean{RC} = \lean{VOA.AffineSl2RigidityCollapse}.

{\footnotesize
\begin{longtable}{>{\raggedright\arraybackslash}p{4.6cm}>{\raggedright\arraybackslash}p{5.0cm}>{\raggedright\arraybackslash}p{2.4cm}>{\raggedright\arraybackslash}p{2.4cm}}
\toprule
Lean declaration & intended meaning & hypotheses & status / check\\
\midrule
\endhead
\lean{Deform.FirstOrderDeformation}, \lean{IsGraded}, \lean{IsCoboundary}, \lean{IsTrivial}, \lean{IsRigidGraded} & first-order deformations, weight rule, coboundaries $\delta\varphi$, graded rigidity & -- & defs.; Def.~\ref{l09:def:rigid}, Sec.~\ref{l09:sec:deform}\\
\lean{Admissible.AdmissibleLevel} & $k+2 = p/q$, $\gcd=1$, $p\ge2$, $q\ge1$ & -- & def.; Def.~\ref{l09:def:admissible}\\
\lean{GR.simple} & $L_k(\lninesl) = V^k/\mathcal{I}_k$ & $k\ne-2$ & def.; Def.~\ref{l09:def:Lk}\\
\lean{GR.simple_isSimple}$^\ast$, \lean{GR.simple_vacuum_ne_zero}$^\ast$ & $L_k$ simple and nonzero & $k\ne-2$ & uncond., pbw-free; Sec.~\ref{l09:sec:affine}\\
\lean{GR.simple_sl2BasisOnWeightOne}$^\ast$ & $V_1(L_k)\cong\lninesl$ & $k\ne0,-2$ & uncond.; Sec.~\ref{l09:sec:affine}\\
\lean{Sl2Whitehead.whitehead2_of_sl2Basis}$^\ast$ & $H^2(\lninesl,\lninesl)=0$ & $\lninesl$-triple basis & uncond.; Thm.~\ref{l09:thm:whitehead}\\
\lean{Deform.weightOneCochain_mem_twoCocycle}$^\ast$ & $a^1_{(0)}b$ is a CE 2-cocycle & graded, $Z(V_1)=0$ & uncond.; Sec.~\ref{l09:sec:whitehead}\\
\lean{LQ.CorrectedSystem.c_eq_zero}$^\ast$ & LQ eqs.\ (10)--(14) force $c=0$ (erratum fixed) & the 5 equations & uncond.; Sec.~\ref{l09:sec:lq}\\
\lean{LQ.lqVector_eq_sl2Singular}$^\ast$, \lean{LQ.maxGradedIdeal_ne_bot}$^\ast$ & LQ's singular vector $=$ library's; $L_{-4/3}\ne V^{-4/3}$ & none & uncond.; Sec.~\ref{l09:sec:lq}\\
\lean{LQ.isRigidGraded_of_split}$^\ast$ & Thm.~3.1 $+$ Piece 1 $+$ Piece 2 $\Rightarrow$ rigid & \lean{hZ,h31,h1,h2} & cond.\ (assembly lemma); Thm.~\ref{l09:thm:split}\\
\lean{LQ.L43_isRigidGraded}$^\ast$ & LQ Thm.~5.1, split form & \lean{hb,h31,h2} (all now proved) & cond., discharged; Thm.~\ref{l09:thm:split}\\
\lean{RC.not_gradedImpliesRegular_L43}$^\ast$, \lean{RC.not_gradedImpliesPoleZero_L43}$^\ast$ & dead end 1: naive \S5 hypothesis is false & none & uncond.; Sec.~\ref{l09:sec:deadends}\\
\lean{RC.not_normalisedPoleVanishes_heisenberg}$^\ast$, \lean{RC.no_uniform_derivation_piece2}$^\ast$, \lean{RC.not_trivial_satisfies_piece2}$^\ast$ & Piece 2 is not tautological (fails for free boson) & $k\ne0$ & uncond.; Sec.~\ref{l09:sec:deadends}\\
\lean{LQ.L43_graded_rigidity}$^\ast$ & \textbf{Linshaw--Qi Thm.~5.1}: $L_{-4/3}$ graded first-order rigid & none & \textbf{uncond.}; Thm.~\ref{l09:thm:L43}\\
\lean{GR.simple_regular_isTrivial}$^\ast$ & regular graded deformations of $L_k$ trivial (Thm.~3.1) & $k\ne0,-2$ & uncond.; Sec.~\ref{l09:sec:uniform}\\
\lean{GR.simple_existsUnique_pole_parameter}$^\ast$ & unique level parameter $\delta$ & graded, pole form & uncond.; Sec.~\ref{l09:sec:uniform}\\
\lean{GR.graded_rigidity_of_relations}$^\ast$ & two relations with $P'(k)\ne0$ $\Rightarrow$ rigid & relations, $P'(k)\neq0$ & uncond.\ criterion; Thm.~\ref{l09:thm:relations}\\
\lean{PolynomialMatrix.kernel_pair_of_simple_root}$^\ast$ & simple det root $\Rightarrow$ kernel pair with $r^{\mathsf T}M's\ne0$ & simple root & uncond.; Lem.~\ref{l09:lem:kernelpair}\\
\lean{GR.graded_rigidity_of_vacuum_simple_root}$^\ast$ & simple root of vacuum Gram det $\Rightarrow$ rigid & $k\ne0,-2$, simple root & uncond.\ criterion; Sec.~\ref{l09:sec:determinant}\\
\lean{vacuumChevalleyDeterminant_admissible_simple}$^\ast$ (ns.\ \lean{Wak.Charged}) & at $k=p/q-2$ the block $(q(p-1),p-1)$ det has a simple root & $p\ge2$, $q\ge1$, coprime & uncond.; Sec.~\ref{l09:sec:determinant}\\
\lean{GR.negative_admissible_graded_rigidity}$^\ast$ & \textbf{$L_k(\lninesl)$ graded first-order rigid for all negative admissible $k$} & $2\le p<2q$, $q\ge1$, coprime & \textbf{uncond.}; Thm.~\ref{l09:thm:headline}\\
\lean{exactPresentation43_of_maximalGeneration}$^\ast$ (and \lean{85}; ns.\ \lean{AffineSl2BoundaryC2Wakimoto}) & $R_{L_k}\cong\C[e,f,h]/(\mathrm{Cas}^n\cdot)$ at $k=-4/3,-8/5$ & maximal generation (Adamovi\'c--Milas) & \textbf{cond.}; Sec.~\ref{l09:sec:boundary}\\
\lean{AdmissiblePartitions.AdmissibleBasisHyp} & Salazar refined character & -- & cited hypothesis, unproved; Sec.~\ref{l09:sec:partitions}\\
\bottomrule
\end{longtable}}

\paragraph{Points a referee should weigh.}
(1) The headline is about the \emph{simple} quotient, defined intrinsically; it is at least as strong as
Linshaw--Qi's notion (Section~\ref{l09:sec:deform}). (2) The ``tautological bundled hypothesis'' of the
library map refers to a \emph{deleted} predicate; no current headline depends on it, and the
remaining conditional split has been fully discharged (Section~\ref{l09:sec:deadends}). (3) The
hypothesis $p<2q$ is used only to get $k\neq0$; the library also states the theorem for all admissible
$k\neq0$ (Theorem~\ref{l09:thm:alladmissible}). (4) $V(c,0)$ is proved not graded rigid, for every $c$ (\lean{not_isRigidGraded_virasoroVOA}); the notion is
non-vacuous (Section~\ref{l09:sec:deform}, second warning). (5)
\lean{RegularImpliesTrivialGraded} is over-general as a predicate on arbitrary $W$ but is only ever
proved at affine $\lninesl$ quotients. (6) Every rigidity statement here is
\emph{graded first-order}; the ungraded \lean{IsRigid} is false for every $V(c,0)$ and is never claimed.

\section{Notes and references}\label{l09:sec:notes}

\begin{sloppypar}
\begin{itemize}[leftmargin=1.5em]
\item A.~R.~Linshaw and F.~Qi, \emph{Deformation rigidity of some simple affine VOAs},
  arXiv:2601.12017 (v2, 2026). Theorem~3.1 (regular $\Rightarrow$ trivial), Theorem~4.1 (positive
  integral levels), Theorem~5.1 ($L_{-4/3}(\lninesl)$), equations (3) and (10)--(14), and the conjecture
  that $L_k(\lnineg)\neq V^k(\lnineg)$ implies rigidity. The erratum in (10) is recorded in
  Section~\ref{l09:sec:lq}.
\item V.~Kovalchuk and F.~Qi, \emph{First-order deformations of freely generated vertex algebras},
  arXiv:2408.16309 --- the deformation complex and the $H^2(\lnineg,\lnineg) = 0$ normalisation (\S6.2).
\item V.~G.~Kac and M.~Wakimoto, \emph{Modular invariant representations of infinite-dimensional Lie
  algebras and superalgebras}, Proc.\ Natl.\ Acad.\ Sci.\ USA \textbf{85} (1988) 4956--4960 --- admissible
  levels and weights.
\item V.~G.~Kac and D.~A.~Kazhdan, \emph{Structure of representations with highest weight of infinite-dimensional
  Lie algebras}, Adv.\ Math.\ \textbf{34} (1979) 97--108 --- the determinant formula and \eqref{l09:eq:KK}.
\item F.~G.~Malikov, B.~L.~Feigin and D.~B.~Fuchs, \emph{Singular vectors in Verma modules over Kac--Moody
  algebras}, Funct.\ Anal.\ Appl.\ \textbf{20} (1986) 103--113.
\item I.~B.~Frenkel and Y.~Zhu, \emph{Vertex operator algebras associated to representations of affine and
  Virasoro algebras}, Duke Math.\ J.\ \textbf{66} (1992) 123--168 --- the affine VOAs $V^k(\lnineg)$ and
  $L_k(\lnineg)$.
\item D.~Adamovi\'c and A.~Milas, \emph{Vertex operator algebras associated to modular invariant representations
  for $A_1^{(1)}$}, Math.\ Res.\ Lett.\ \textbf{2} (1995) 563--575 --- the maximal ideal at admissible
  level is generated by the singular vector (their Theorem~3.1.4, the external input of
  Section~\ref{l09:sec:boundary}).
\item D.~Adamovi\'c, \emph{A construction of admissible $A_1^{(1)}$-modules of level $-4/3$}, J.\ Pure Appl.\
  Algebra \textbf{196} (2005) 119--134 --- the weight-3 singular vector \eqref{l09:eq:sing43}.
\item H.~Li, \emph{Symmetric invariant bilinear forms on vertex operator algebras}, J.\ Pure Appl.\ Algebra
  \textbf{96} (1994) 279--297 --- the radical of the invariant form is the maximal ideal.
\item B.~H.~Lian, \emph{On the classification of simple vertex operator algebras}, Comm.\ Math.\ Phys.\
  \textbf{163} (1994) 307--357 --- cited by Linshaw--Qi for $W_{(0)}=\C\lnineone$; the library avoids it.
\item M.~Gorelik and V.~G.~Kac, \emph{On simplicity of vacuum modules}, Adv.\ Math.\ \textbf{211} (2007)
  621--677 --- criterion for $V^k(\lnineg)$ to be simple (used only in Remark~\ref{l09:rem:extension}).
\item T.~Arakawa, \emph{Rationality of admissible affine vertex algebras in the category $\mathcal{O}$},
  Duke Math.\ J.\ \textbf{165} (2016) 67--93, arXiv:1207.4857; T.~Arakawa, C.~Jiang and A.~Moreau,
  on the $C_2$ algebras of affine VOAs at boundary admissible levels (the symbol nonvanishing used in
  \lean{ClassicalData}).
\item D.~Salazar, arXiv:2405.07126v3 (refined characters of the boundary Virasoro minimal models via Andrews--Gordon) --- the input
  named \lean{AdmissibleBasisHyp}.
\item C.~Beem, M.~Lemos, P.~Liendo, W.~Peelaers, L.~Rastelli and B.~C.~van Rees, \emph{Infinite chiral
  symmetry in four dimensions}, Comm.\ Math.\ Phys.\ \textbf{336} (2015) 1359--1433, arXiv:1312.5344 ---
  chiral algebras of 4d $\mathcal{N}=2$ theories; T.~Creutzig, \emph{W-algebras for Argyres--Douglas
  theories}, arXiv:1701.05926 --- boundary admissible $L_k(\lninesl)$ and $(A_1,D_{2n+1})$ theories.
\item A.~M.~Polyakov, \emph{Quantum gravity in two dimensions}, Mod.\ Phys.\ Lett.\ A \textbf{2} (1987)
  893 --- fractional-level $SL(2)$ current algebra in 2d gravity.
\item Library sources: \leanpath{VOA/Examples/AffineSl2Rigidity43.lean} (split, erratum, dead ends),
  \leanpath{VOA/Examples/AffineSl2RigidityCollapse.lean} (refutation certificates),
  \leanpath{VOA/Core/WeightOneCohomology.lean} (Whitehead), \leanpath{VOA/Examples/AffineSl2General*.lean}
  and \leanpath{AffineSl2UniformObstruction.lean} (uniform machinery),
  \leanpath{VOA/Examples/AffineSl2Admissible*.lean} (admissible arithmetic, simple root, headline);
  the book's catalogue of dead ends, Section~\ref{ch15:sec:deadends}.
\end{itemize}
\end{sloppypar}
